%!TEX program = xelatex
%!TEX encoding = UTF-8
% ============================================================================
%  功率控制算法与超级电容技术文档
%  模板：ElegantNote (https://github.com/ElegantLaTeX/ElegantNote)
%  编译：XeLaTeX + Biber
% ============================================================================
\documentclass[cn,blue,11pt,device=normal]{elegantnote}

%------------------------------------------------------------------ 页面
% device=normal  ->  A4 纸张，1 英寸页边距
% 模板的 geye / hazy / sepia 模式会给纸张上底色，这里显式恢复为纯白背景。
\definecolor{paperwhite}{RGB}{255,255,255}
\pagecolor{paperwhite}

%------------------------------------------------------------------ 宏包
\usepackage{array}
\usepackage{amsmath}
\usepackage{amssymb}
\usepackage{booktabs}
\usepackage{multirow}
\usepackage{pgfplots}
\pgfplotsset{compat=1.17}
\usetikzlibrary{arrows.meta,positioning,calc,fit,backgrounds,shapes.geometric,patterns}
\usepackage{algorithm}
\usepackage{algpseudocode}
\usepackage{listings}

%------------------------------------------------------------------ 版式微调
% A4 正文较宽，适当收紧行距与段距，避免过于松散。
\linespread{1.28}
\setlength{\parskip}{2.2mm}
\setlength{\abovedisplayskip}{5pt}
\setlength{\belowdisplayskip}{5pt}
\setlength{\abovedisplayshortskip}{3pt}
\setlength{\belowdisplayshortskip}{3pt}

% 中英混排时给 TeX 一点额外伸缩余量，显著减少 overfull \hbox
\setlength{\emergencystretch}{2em}
\tolerance=2000
\hbadness=2500

% 参考文献中的长 URL 允许在数字/大写/小写字符间断行
\setcounter{biburlnumpenalty}{100}
\setcounter{biburlucpenalty}{100}
\setcounter{biburllcpenalty}{100}

% 一级标题（部分）使用模板主题色，居中显示
\titleformat{\part}[display]
  {\centering\normalfont\huge\bfseries\color{ecolor}}
  {\Large\color{ecolor}Part \Roman{part}}
  {1.4em}
  {\Huge}
\renewcommand{\thepart}{\Roman{part}}
\titlespacing*{\part}{0pt}{2\baselineskip}{2\baselineskip}

% 表注 / 图注字号
\captionsetup{font=small}
\captionsetup[table]{position=top}

%------------------------------------------------------------------ 算法环境汉化
\floatname{algorithm}{算法}
\renewcommand{\algorithmicrequire}{\textbf{输入：}}
\renewcommand{\algorithmicensure}{\textbf{输出：}}
\renewcommand{\algorithmicreturn}{\textbf{返回}}
\renewcommand{\algorithmicfor}{\textbf{for}}
\renewcommand{\algorithmicdo}{\textbf{do}}
\renewcommand{\algorithmicwhile}{\textbf{while}}
\renewcommand{\algorithmicif}{\textbf{if}}
\renewcommand{\algorithmicthen}{\textbf{then}}
\renewcommand{\algorithmicelse}{\textbf{else}}
\renewcommand{\algorithmicend}{\textbf{end}}

%------------------------------------------------------------------ 代码样式
\definecolor{codegray}{RGB}{245,246,248}
\definecolor{codeframe}{RGB}{214,219,226}
\definecolor{codegreen}{RGB}{0,120,80}
\definecolor{codepurple}{RGB}{150,40,120}

\lstdefinestyle{cstyle}{
  language=C,
  basicstyle=\ttfamily\footnotesize,
  keywordstyle=\color{eblue}\bfseries,
  commentstyle=\color{codegreen}\itshape,
  stringstyle=\color{codepurple},
  numberstyle=\tiny\color{gray},
  numbers=left,
  numbersep=9pt,
  frame=single,
  framesep=5pt,
  rulecolor=\color{codeframe},
  backgroundcolor=\color{codegray},
  breaklines=true,
  breakatwhitespace=false,
  postbreak=\mbox{\textcolor{gray}{$\hookrightarrow$}\space},
  tabsize=4,
  showstringspaces=false,
  columns=fullflexible,
  keepspaces=true,
  upquote=true,
  morekeywords={fp32,uint8_t,uint16_t,int32_t,chassis_move_t,cap_measure_t,RC_ctrl_t},
  literate={　}{{\hspace{1em}}}1
}

\lstdefinestyle{matstyle}{
  language=Matlab,
  basicstyle=\ttfamily\footnotesize,
  keywordstyle=\color{eblue}\bfseries,
  commentstyle=\color{codegreen}\itshape,
  stringstyle=\color{codepurple},
  numberstyle=\tiny\color{gray},
  numbers=left,
  numbersep=9pt,
  frame=single,
  framesep=5pt,
  rulecolor=\color{codeframe},
  backgroundcolor=\color{codegray},
  breaklines=true,
  tabsize=4,
  showstringspaces=false,
  columns=fullflexible,
  keepspaces=true
}

%------------------------------------------------------------------ 框图统一样式
% 全文所有功能框图共用这一套样式，保证风格一致。
\tikzset{
  fblock/.style   = {draw=eblue, thick, rounded corners=2pt, fill=eblue!5,
                     minimum height=9mm, minimum width=20mm, align=center,
                     font=\small, inner sep=3pt},
  fsys/.style     = {fblock, draw=codegreen, fill=codegreen!7, minimum width=26mm},
  fwarn/.style    = {fblock, draw=red!65!black, fill=red!5},
  fsoft/.style    = {fblock, fill=eblue!10, minimum width=24mm},
  fsum/.style     = {draw=eblue, thick, circle, inner sep=0pt,
                     minimum size=5.5mm, fill=white},
  fio/.style      = {draw=gray!65, rounded corners=7pt, fill=gray!6,
                     align=center, font=\small, inner sep=4pt},
  fgroup/.style   = {draw=gray!55, dashed, rounded corners=3pt, inner sep=7pt},
  farr/.style     = {draw, thick, -{Stealth[length=1.8mm]}},
  fdarr/.style    = {draw, thick, dashed, -{Stealth[length=1.8mm]}},
  fbio/.style     = {draw, thick, {Stealth[length=1.8mm]}-{Stealth[length=1.8mm]}},
  flab/.style     = {font=\scriptsize, text=black!70, inner sep=1pt},
  ftitle/.style   = {font=\small\bfseries, text=eblue}
}

% 行内代码：等宽、可在空格与下划线处断行。
% 注意不要用 \mbox 包住——那会变成一个不可断行的盒子，
% 长标识符（如 motor_speed_pid[i].out）会直接冲出页边。
\newcommand{\code}[1]{{\ttfamily\small\def\_{\textunderscore\allowbreak}#1}}

%------------------------------------------------------------------ 记号宏
% 注意每个宏末尾都带一个空的 {}：这些宏本身已经含下标（如 P_{\mathrm{in}}），
% 若在宏名后面直接跟一个下标，会形成“同一个基底下标两次”，
% 触发 Double subscript 错误，第二个下标被 TeX 静默丢弃（页面上看不出来）。
% 末尾的 {} 提供一个零宽基底，因此“宏名 + 下标”与
% “用大括号包住宏名再加下标”两种写法都能正常工作。
\newcommand{\Pin}{P_{\mathrm{in}}{}}
\newcommand{\Pmech}{P_{\mathrm{m}}{}}
\newcommand{\Icmd}{I_{\mathrm{cmd}}{}}
\newcommand{\Itau}{I_{\tau}{}}
\newcommand{\KT}{K_{T}{}}
\newcommand{\CT}{C_{T}{}}
\newcommand{\Ebuf}{E_{\mathrm{buf}}{}}
\newcommand{\Plim}{P_{\mathrm{lim}}{}}
\newcommand{\Pmax}{P_{\max}{}}
\newcommand{\dd}{\mathrm{d}}
\newcommand{\fref}[1]{式~\eqref{#1}}

%------------------------------------------------------------------ 其他
\newcommand{\pending}[1]{%
  \par\medskip\noindent
  \begin{note}#1\end{note}\par\medskip}

\graphicspath{{image/}{images/}{figures/}}

% ============================================================================
\title{功率控制算法与超级电容技术文档}

\author{TuxMonkey}
\institute{\href{https://github.com/TuxMonkey2004/Docs-of-PowerControl-and-SuperCap}{Docs-of-PowerControl-and-SuperCap}}

\version{1.00}
\date{\zhdate{2026/10/05}}

\begin{document}

\maketitle

\centerline{\includegraphics[width=0.18\textwidth]{logo-blue.png}}

%------------------------------------------------------------------ 摘要
\phantomsection
\section*{摘要}
\addcontentsline{toc}{section}{摘要}

本文围绕 RoboMaster 超级电容（下称“超电”）车种的底盘功率控制问题展开。超电接入之后，
传统的“缓冲能量闭环 + 力矩阈值”方案失去了可观测性：理想超电会无条件补偿超出功率上限的能量，
缓冲能量不再反映电机真实需求，因而必须直接估计电机的输入功率。本文以 M3508 电机与 C620 电调为对象，
完整推导了由控制量到力矩、再由力矩到输入功率的物理链路，建立了
$\Pin = k_1\omega^2 + \CT\,\Icmd\,\omega + k_2\Icmd^2 + a$ 形式的电机功率模型，
补齐了原资料中略去的\textbf{最小二乘参数辨识}的完整数学推导（含准则函数、一阶必要条件、
正规方程、可解性条件、量纲与数值条件、加权与鲁棒回归的修正方法），
并对“给定功率反解力矩”这一步做了逐项分析：判别式的含义、两根的物理来源（回馈制动）、
根的选择规则，以及原实现中未覆盖的两类边界情形及其修复方案。

在此基础上，本文比较了均分功率、绝对值统一缩放与本文采用的\textbf{功率再分配}
三种多电机分配策略，给出了各自的误差来源分析；随后分析了缓冲能量闭环的结构与
参数整定方法，指出该闭环的被控对象是纯积分器，并给出闭环特征方程
$s^2 + K_p s + K_i = 0$ 与阻尼比、自然频率的对应关系。最后逐行解读了开源实现，
列出了其中若干处可能导致行为偏差的实现细节，并给出改进建议。

全文分为两个部分：\textbf{Part I} 为功率控制，\textbf{Part II} 为超级电容，
后者目前仅给出总体框架，详细内容待后续补充。

\keywords{底盘功率控制；无刷电机功率模型；最小二乘参数辨识；功率再分配；缓冲能量闭环；超级电容}
\medskip

\newpage
\phantomsection
\tableofcontents
\newpage

% ============================================================================
\part{功率控制}

% ============================================================================
\section{背景、问题与符号约定}

\subsection{缓冲能量是怎么工作的}

裁判系统发给底盘的并不是一个“绝对不能超过”的功率上限，而是把“还能多用多少”
和“已经用超了多少”合在同一个量里管理。这个量就是\emph{缓冲能量}
（buffer energy）$\Ebuf$，它的变化规律只有一条：
\begin{equation}\label{eq:buffer}
  \frac{\dd \Ebuf}{\dd t} = \Pmax - \Pin ,
\end{equation}
其中 $\Pmax$ 为当前等级对应的底盘功率上限，$\Pin$ 为底盘从电池实际取用的功率。
\fref{eq:buffer} 的含义非常直白：

\begin{itemize}
  \item $\Pin < \Pmax$：底盘用得比上限少，差额先存下来，$\Ebuf$ 变大；
  \item $\Pin > \Pmax$：底盘用得比上限多，超出的部分从这里扣，$\Ebuf$ 变小。
\end{itemize}

也就是说，$\Ebuf$ 就是\emph{功率差对时间的累积}，单位是焦耳。
它有一个上限 $E_{\max}$（常见取值为 $60\ \mathrm{J}$，随赛季与兵种等级可能不同）：
到达上限后就不再增加；一旦被扣到 $0$ 之后还在超功率，裁判系统就开始扣血\cite{rm2024rules}。
整个过程如图~\ref{fig:bufferbehaviour} 所示。

\begin{figure}[htbp]
  \centering
  \begin{tikzpicture}[x=1.0cm, y=0.048cm, >=Stealth]
    % 扣血区域
    \fill[red!10] (8.2,0) rectangle (11.0,64);
    % 坐标轴
    \draw[gray!70, thick, -{Stealth[length=1.8mm]}] (0,0) -- (11.5,0)
      node[right, font=\small] {$t$};
    \draw[gray!70, thick, -{Stealth[length=1.8mm]}] (0,0) -- (0,68)
      node[above, font=\small] {$\Ebuf$ (J)};
    % 60 J 上限
    \draw[red!60!black, dashed] (0,60) -- (11.0,60);
    \node[flab, text=red!60!black, anchor=south west] at (0.0,61) {$E_{\max}=60\ \mathrm{J}$};
    % 缓冲能量曲线
    \draw[eblue, very thick]
      (0,18) -- (1.6,44) -- (3.4,14) -- (5.2,60) -- (6.4,60) -- (8.2,0) -- (11.0,0);
    % 标注（白底，避免文字与曲线互相干扰）
    \node[flab, fill=white, inner sep=1.5pt] at (0.55,40) {上升};
    \node[flab, fill=white, inner sep=1.5pt] at (2.45,27) {下降};
    \node[flab, anchor=south] at (5.80,62)
      {到达上限后不再增加};
    \node[flab, text=red!65!black, align=center, anchor=north] at (9.20,55)
      {缓冲能量耗尽\\开始扣血};
    \node[flab, anchor=west] at (0.10,10) {$0$};
  \end{tikzpicture}
  \caption{缓冲能量的典型变化过程。它本身不是一个可以“取用”的储量，
  而是功率差随时间的累积：底盘功率低于上限 $\Pmax$ 时上升，超过 $\Pmax$ 时下降；
  到达 $60\ \mathrm{J}$ 后饱和不再增加，被扣到 $0$ 之后进入扣血阶段。}
  \label{fig:bufferbehaviour}
\end{figure}

由此可以得到一个关键结论：\textbf{缓冲能量是判断底盘有没有超功率的唯一观测量}。
只要底盘长期工作在 $\Pmax$ 附近，$\Ebuf$ 就会在某个值附近来回摆动。
过去几年的功率控制方案，本质上都是围绕 $\Ebuf$ 做一个外环：

\begin{quote}
用 PID 把 $\Ebuf$ 稳定在某个设定值（例如 $30\ \mathrm{J}$）：
$\Ebuf$ 偏低，说明前段时间用得太多，就把本周期的功率限额收紧一点；
$\Ebuf$ 偏高，说明还有余量，就放宽一点。
这样底盘既不会长期超功率，又能在需要的时候短时间超过 $\Pmax$ 冲一下。
\end{quote}

\subsection{超电为什么会让缓冲能量失效}

超级电容通过双向 DC--DC 变换器并联在电池母线与电调输入端之间，其行为可以概括为：
\begin{itemize}
  \item 底盘需求功率 $\Pin^{\mathrm{req}} > \Pmax$ 时，电容放电补足差额，母线电压维持不变；
  \item 底盘需求功率 $\Pin^{\mathrm{req}} < \Pmax$ 时，多余的能量被用于给电容充电。
\end{itemize}

一个“理想”的超电会\textbf{无条件}地把差额补齐。于是从裁判系统的视角看，
底盘的输入功率被削平在 $\Pmax$ 附近，而这个“削平”与四个电机到底想要多少功率已经无关了。
\fref{eq:buffer} 因此退化为 $\dd \Ebuf/\dd t \approx 0$：缓冲能量不再随电机需求变化，
\begin{equation}
  \frac{\partial \Ebuf}{\partial \Pin^{\mathrm{req}}} \approx 0 .
\end{equation}
控制理论的语言是：被控量与控制量之间的通道被超电“短路”了，外环失去了对电机的可观性。
继续用缓冲能量闭环去限制电机功率，只会得到“无论怎么调 PID 都限不住”的结果，
而且会持续掏空电容电量。

结论是明确的：一旦引入超电，就必须\textbf{直接获得电机的功率信息}，而不能再依赖缓冲能量。

\subsection{两条技术路线}

获得电机功率只有两条路。

\begin{description}
  \item[路线 A：直接测量。] 在每个电调输入端串入电流采样（如 INA226 配合分流电阻），
  测量母线电流与电压。优点是无需建模、结果直接；缺点是硬件复杂度与成本上升，
  采样链路（分流电阻、隔离、ADC 带宽、共地）引入新的失效点，且需要为四个电机各配一路。
  更麻烦的是，电流波形含大量开关纹波，要得到“平均功率”还需要合理的滤波，
  滤波又会带来相位滞后，反过来恶化控制品质。
  \item[路线 B：模型估计。] 利用电调已经回传的两个量——\emph{力矩电流控制值} $\Icmd$
  与\emph{转子转速} $\omega$——通过一个电机功率模型估算 $\Pin$。优点是不增加任何硬件，
  复用已有反馈；缺点是完全依赖模型的准确性。
\end{description}

本文采用的是路线 B（官方例程中的功率控制实现可参考文献~\cite{rmdevboard}）。它成立的前提是一个可以被验证的假设：\textbf{C620 电调的力矩电流环
足够快}，以至于“发送的电流指令”与“电机实际流过的力矩电流”在同一控制周期内是等价的。
原作者通过对比电调回传的 \code{given\_current} 与发送的 \code{give\_current} 印证了这一假设，
我们将在第~\ref{sec:cmd2torque} 节中给出这一环节的完整量纲推导。

\subsection{问题的本质：从“力矩约束”到“功率约束”}

在得到模型之前，工程上最朴素的做法是限制总力矩：把四个电机的目标力矩求和，
若超过某个预设阈值，就等比缩小。这个做法有一个结构性缺陷。
电机的输入功率近似为
\begin{equation}
  \Pin \approx \frac{\tau\,\omega}{9.55} + P_{\mathrm{loss}}(\tau,\omega),
\end{equation}
它\emph{同时}依赖力矩 $\tau$ 与转速 $\omega$。因此：

\begin{itemize}
  \item \textbf{低速时}，同一个力矩对应的功率很小。用固定的力矩阈值去限制，
  相当于在低速段给了一个“过紧”的约束，底盘起步疲软、加速缓慢；
  \item \textbf{高速时}，同一个力矩对应的功率很大，力矩阈值反而\emph{过松}，
  会造成严重超功率。
\end{itemize}

换言之，力矩阈值与功率上限之间不存在一个与工况无关的换算关系。
正确的做法是反过来问：

\begin{quote}
\textbf{在当前转速下，为了恰好消耗给定功率 $P^\ast$，我应该发送多大的力矩指令？}
\end{quote}

把这个问题回答清楚，就能实现一个理想的功率控制目标——\emph{恒功率运行}：
从起步到大加速度的整个过程中，电机始终工作在 $\Pmax$ 上，从而获得功率受限条件下
理论上的最大加速度。这一目标也解释了为什么必须建模：它要求的是“功率到力矩”的\emph{反函数}。

\subsection{为什么必须“预测”功率}\label{sec:whypredict}

既然缓冲能量能反映超功率，为什么不能“先用着，超了再收”？
这个思路在底盘上会遇到三个绕不过去的困难。

\paragraph{困难一：等测到，就已经晚了。}
裁判系统对底盘功率的结算周期是 $100\ \mathrm{ms}$（即 $10\ \mathrm{Hz}$）\cite{rm2024rules}，
而底盘的控制周期是 $1\ \mathrm{ms}$ 量级。也就是说，
从“电机真的超功率了”到“程序知道超功率了”，中间隔着一百个控制周期。
在这段时间里底盘只能继续超功率，等数据回来再收手，多消耗的能量已经发生了。
这是一条“先犯错、再纠正”的路，而功率约束恰恰不给人犯错的机会：
缓冲能量被扣到 $0$ 之后，每一次超功率都直接扣血。

\paragraph{困难二：反馈量分不到具体电机。}
裁判系统给出的只有\emph{四个电机的总功率}。当总功率超标时，
程序并不知道是哪一个电机超了——可能只有右前轮在打滑空转，
也可能四个轮子一起用力过猛。而功率限制最终必须落到每一个电机的力矩指令上，
这中间缺了一个“总功率 $\to$ 各电机力矩”的分配环节，而反馈信息里没有它。

\paragraph{困难三：能下发的量不是功率。}
这是最根本的一点。程序能直接控制的量是\emph{力矩电流指令} $\Icmd$，
而功率与指令之间不是固定的倍数关系：
\begin{equation}\label{eq:why-nonlinear}
  \Pin = k_2\omega^2 + \CT\,\Icmd\,\omega + k_1\Icmd^2 + a .
\end{equation}
同一个指令 $\Icmd$，在 $1000\ \mathrm{rpm}$ 下可能只消耗几瓦，
在 $6000\ \mathrm{rpm}$ 下却要消耗几十瓦；反过来，同样要求消耗 $40\ \mathrm{W}$，
低速时要发很大的指令，高速时只需要很小的指令。
所以“功率超了就把指令调小一点”这种做法，\emph{调多少完全取决于当前转速}，
并不存在一个通用的比例。

\medskip
三个困难指向同一个结论：\textbf{与其在下发之后去测量和纠正，
不如在下发之前就把结果算出来}。这就是“预测”的含义：

\begin{quote}
用电机功率模型，从\emph{准备下发的指令}与\emph{当前转速}，
直接算出\emph{这样下发会消耗多少功率}；反过来，
也可以从\emph{希望消耗多少功率}算出\emph{应该下发多大的指令}。
\end{quote}

两种思路的差别如图~\ref{fig:why-predict} 所示。反馈式控制必须“先发生、再发现”，
而预测式控制在下发之前就已经知道结果。

\begin{figure}[htbp]
  \centering
  \begin{tikzpicture}[>=Stealth]
    % ================= (a) 反馈式 =================
    \node[ftitle] at (-1.35,0) {(a)};
    \node[fblock] (a1) at (0,0)    {力矩指令};
    \node[fblock] (a2) at (3.4,0)  {电机};
    \node[fblock] (a3) at (6.8,0)  {实际功率};
    \node[fblock] (a4) at (10.2,0) {与上限比较};
    \draw[farr] (a1) -- (a2);
    \draw[farr] (a2) -- (a3);
    \draw[farr] (a3) -- (a4);
    \draw[fdarr] (a4.south) -- ++(0,-0.9) -| (a1.south);
    \node[flab, text=red!65!black, anchor=north] at (5.1,-1.75)
      {要等下一个结算周期（$100\ \mathrm{ms}$）才生效——此刻已经超功率了};

    % ================= (b) 预测式 =================
    \node[ftitle] at (-1.35,-4.1) {(b)};
    \node[fblock] (b1) at (0,-4.1)    {目标功率\\$P^{\ast}$};
    \node[fsoft]  (b2) at (3.4,-4.1)  {模型反解\\$f^{-1}$};
    \node[fblock] (b3) at (6.8,-4.1)  {力矩指令\\$\Icmd$};
    \node[fblock] (b4) at (10.2,-4.1) {电机};
    \draw[farr] (b1) -- (b2);
    \draw[farr] (b2) -- (b3);
    \draw[farr] (b3) -- (b4);

    \node[fsoft] (b5) at (3.4,-6.5) {功率预测\\$\hat\Pin = f(\Icmd,\omega)$};
    \node[flab, anchor=east] at (1.75,-6.5) {$\Icmd$、$\omega$};
    \draw[farr] (1.85,-6.5) -- (b5.west);
    \draw[fdarr] (b5.north) -- (b2.south);
    \node[flab, text=codegreen, anchor=north west] at (2.0,-7.35)
      {同一套模型正反两用：下发前先算一遍，因此不需要“事后修正”};
  \end{tikzpicture}
  \caption{两种控制思路的对比。(a) 反馈式：先下发、再测量、再修正，
  修正永远滞后一个结算周期；(b) 预测式：用同一个模型在下发之前算出后果，
  并把“想要多少功率”直接翻译成“该发多大指令”。}
  \label{fig:why-predict}
\end{figure}

预测带来三个直接的好处：
\begin{enumerate}
  \item 功率是在下发\emph{之前}分配好的，不存在“先超标再收回”的窗口；
  \item 预算可以精确地分给四个电机，而不是靠一个笼统的系数去凑；
  \item \emph{恒功率运行}成为可能——从起步到加速的整个过程，电机始终贴着功率上限走。
  这正是功率受限条件下加速度最大的工作方式（参见图~\ref{fig:envelope}）。
\end{enumerate}

当然，预测不是免费的：它的准确性完全取决于模型。
这就是第~\ref{sec:model} 节与第~\ref{sec:lsq} 节要花大篇幅建立模型、
做参数辨识的原因。而模型误差又不可能完全消除（温度、母线电压、电机个体差异都会带来偏差），
所以除了预测，还必须保留第~\ref{sec:bufferloop} 节的缓冲能量外环作为兜底——
一句话概括两者的分工：\textbf{预测负责“准”，外环负责“稳”}。

\subsection{系统总体结构}\label{sec:system}

把上面的思路落到实车上，就得到一个分层的系统结构，如图~\ref{fig:system} 所示。
从上到下可以分成四层。

\begin{description}
  \item[感知层：] 提供全部输入——操作手的遥控指令、
  裁判系统给出的功率上限与缓冲能量、以及四个电机回传的转速与力矩电流。
  \item[决策层（主控软件）：] 全部算法都在这一层，而且分成两个节奏。
  外环按裁判系统的 $10\ \mathrm{Hz}$ 结算节奏运行，负责把缓冲能量稳定住并给出功率预算；
  内环按 $1\ \mathrm{kHz}$ 运行，负责预测功率、分配预算、反解指令。
  \item[执行层：] 通过 CAN 把力矩指令发给 C620 电调，由电调驱动 M3508 电机。
  \item[能量层：] 电池母线与超级电容之间通过双向 DC--DC 交换能量，
  在底盘需要爆发时补足差额，在底盘有余量时把能量收回去。
\end{description}

\begin{figure}[htbp]
  \centering
  \begin{tikzpicture}[>=Stealth, transform shape, scale=0.94]
    % ---------------- 感知层 ----------------
    \node[fio] (rc)  at (0,0)    {操作手\\遥控指令};
    \node[fio] (ref) at (5.0,0)  {裁判系统\\$\Pmax$、$\Ebuf$};
    \node[fio] (fb)  at (10.0,0) {电机反馈\\$\omega$、$\Icmd$};

    % ---------------- 外环 ----------------
    \node[fblock] (bloop) at (5.0,-2.6) {缓冲能量\\外环 PID};
    \node[fblock] (bud)   at (10.0,-2.6){功率预算\\$\Plim$};
    \draw[farr] (ref) -- (bloop);
    \draw[farr] (bloop) -- (bud);

    % ---------------- 内环 ----------------
    \node[fsoft]  (pred)  at (0,-5.2)    {功率预测\\$\hat P_i=f({\Icmd}_i,\omega_i)$};
    \node[fblock] (alloc) at (5.0,-5.2)  {功率再分配\\$k=\Plim/\sum_i \hat P_i$};
    \node[fblock] (inv)   at (10.0,-5.2) {反解\\${\Icmd}_i'=f^{-1}(P_i^{\ast},\omega_i)$};
    \draw[farr] (rc) -- (pred);
    \draw[farr] (pred) -- (alloc);
    \draw[farr] (alloc) -- (inv);
    \draw[farr] (bud) -- (alloc);

    % ---------------- 执行层 ----------------
    \node[fblock] (can) at (0,-7.8)    {CAN 发送};
    \node[fblock] (esc) at (5.0,-7.8)  {C620 电调};
    \node[fblock] (mot) at (10.0,-7.8) {M3508 电机 $\times4$};
    \draw[farr] (inv.south) -- ++(0,-1.0) -| (can.north);
    \draw[farr] (can) -- (esc);
    \draw[farr] (esc) -- (mot);

    % ---------------- 能量层 ----------------
    \node[fio]  (bat) at (0,-10.4)   {电池母线};
    \node[fsys] (cap) at (5.0,-10.4) {超级电容\\$+$ 双向 DC--DC};
    \node[flab, anchor=west] at (6.9,-10.4) {补足差额 / 收回余量};
    \draw[fbio] (bat) -- (cap);

    % ---------------- 反馈回路 ----------------
    \draw[fdarr] (mot.east) -- ++(1.5,0) |- (fb.east);
    \node[flab, anchor=west, text=codegreen] at (11.6,-3.9) {反馈};

    % ---------------- 软件范围 ----------------
    \node[fgroup, fit=(bloop)(bud)(pred)(alloc)(inv),
          label={[flab, anchor=south west]above left:主控软件（决策层）}] {};
  \end{tikzpicture}
  \caption{底盘功率控制系统的总体结构。虚线框内是运行在主控上的软件：
  外环按裁判系统的结算节奏（$10\ \mathrm{Hz}$）给出功率预算，
  内环按控制周期（$1\ \mathrm{kHz}$）完成功率预测、预算分配与指令反解；
  超级电容在能量层与电池母线并联，负责补足或收回差额。}
  \label{fig:system}
\end{figure}

这个结构里有两条信息流值得单独指出。

\begin{itemize}
  \item \textbf{纵向：预算的传递。} 缓冲能量外环不直接参与电机运算，
  它只输出一个数——本周期的功率预算 $\Plim$。内环把它当作“总共能用多少功率”，
  再决定怎么分给四个电机。外环与内环解耦，是这套结构最容易维护的地方：
  调外环参数不会影响内环的分配逻辑，反之亦然。
  \item \textbf{横向：模型的复用。} 内环里“功率预测”与“反解”用的是\emph{同一个模型}，
  只是一个正向用、一个求反。因此模型参数只需要标定一次，
  任何一处改进（例如加入母线电压修正）都会同时改善预测与反解。
\end{itemize}

\subsection{符号与正方向约定}\label{sec:convention}

全文严格遵守下表约定，并且\textbf{在整个推导过程中不使用绝对值}。
这一点看似琐碎，实际上是理解第~\ref{sec:inverse} 节“两个根”问题的前提。

\begin{table}[htbp]
  \centering
  \small
  \caption{符号与正方向约定}
  \label{tab:symbols}
  \begin{tabular}{@{}lll@{}}
    \toprule
    符号 & 含义 & 约定 \\
    \midrule
    $\tau$          & 电机转子输出力矩（N$\cdot$m） & 顺时针为正，逆时针为负 \\
    $\omega$        & 电机转子转速（rpm）           & 顺时针为正，逆时针为负 \\
    $\Pin$          & 电机输入功率（W）             & 消耗功率为正，回馈功率为负 \\
    $\Pmech$        & 机械功率（W）                 & $= \tau\omega/9.55$ \\
    $\Icmd$         & 力矩电流控制值（发送值）      & 无量纲，范围 $[-16384,\,16384]$ \\
    $\Itau$         & 力矩电流（A）                 & $= \tfrac{20}{16384}\Icmd$ \\
    $\KT$           & 转子侧力矩常数（N$\cdot$m/A） & $= 0.01562$ \\
    $\CT$           & 力矩电流系数                  & $= 1.9969\times 10^{-6}$ \\
    $\Ebuf$         & 缓冲能量（J）                 & 裁判系统观测量 \\
    $\Pmax$         & 底盘功率上限（W）             & 裁判系统给定 \\
    \bottomrule
  \end{tabular}
\end{table}

\begin{remark}
“正功率表示消耗、负功率表示输出”这一条约定，在第~\ref{sec:redistribute} 节讨论
多电机功率再分配时起着决定性作用：当一个电机处于回馈制动状态时，
它的功率是\emph{负}的，可以把其他电机的功率预算“腾”出来。
任何在建模阶段取绝对值的做法，都会永久性地丢失这一部分信息。
\end{remark}

% ============================================================================
\section{从控制量到力矩：C620 与 M3508 的物理链路}\label{sec:cmd2torque}

这一节回答一个困扰初学者的具体问题：\emph{我发下去的那个数字，究竟代表什么？}

\subsection{三个容易混淆的“电流”}

\begin{description}
  \item[Give current（发送值）：] 主控通过 CAN 报文发给 C620 电调的\emph{控制量} $\Icmd$，
  取值范围 $-16384 \sim 0 \sim 16384$。它是一个无量纲的整型数，物理含义是
  “目标力矩电流”的归一化表示。
  \item[Given current（回馈值）：] C620 在状态报文里回传的\emph{反馈量}，量纲与取值范围
  和发送值相同。它反映了电调内部电流环的实际跟踪情况。两者在稳态下应当基本重合；
  如果长期偏离，说明力矩电流环带宽不足，路线 B 的前提假设不成立。
  \item[真实母线电流：] 电调从母线抽取的实际电流，需要用 INA226 之类的采样电路测量。
  \textbf{它与前两者不是一回事}，不能用它乘母线电压去反推力矩。
\end{description}

\subsection{力矩电流不是相电流}

一个非常常见的误解是“力矩电流乘母线电压就是功率”。事实并非如此。

无刷电机采用磁场定向控制（FOC）时，定子电流被分解为 $d$ 轴分量 $i_d$ 与 $q$ 轴分量 $i_q$。
在 $i_d = 0$ 控制下，$i_q$ 与电磁转矩成正比，而\emph{相电流}是三相绕组中实际流过的
瞬时电流，其幅值与 $i_q$ 相关但相位随转子位置旋转；\emph{母线电流}则还要经过逆变桥的
开关调制，与相电流之间差一个与调制度、功率因数有关的比例因子。

因此：
\begin{itemize}
  \item 力矩电流 $\Itau \leftrightarrow i_q$：与转矩成正比，可以线性换算；
  \item $\Itau \times U_{\mathrm{bus}}$ \textbf{没有}功率的物理意义。
\end{itemize}

真正把两者联系起来的是\emph{损耗}：铜损正比于电流的平方，这正是后文模型中
$k_1\Icmd^2$ 项的来源（见第~\ref{sec:loss} 节）。

\subsection{力矩电流的量纲换算}

C620 的力矩电流控制值满量程 $\pm 16384$ 对应输出力矩电流 $\pm 20\ \mathrm{A}$（手册~\cite{c620manual}），
因此
\begin{equation}\label{eq:itau}
  \Itau = \frac{20}{16384}\,\Icmd \qquad [\mathrm{A}].
\end{equation}
这是一个纯粹的线性标度变换，使用时要注意 $20/16384 = 1.2207\times 10^{-3}$ 这一系数
在整条链路中会反复出现。

\subsection{力矩常数与减速比}

电机手册~\cite{m3508manual} 给出的转矩常数是以\emph{减速箱输出轴}为基准的常数
$K_T^{\mathrm{out}} = 0.3\ \mathrm{N\cdot m/A}$。而 C620 回传的转速是\emph{转子}转速，
力矩也必须统一到转子侧，所以需要按减速比折算。

M3508 搭配的 P19 行星减速箱减速比为
\begin{equation}
  i = \frac{3591}{187} \approx 19.203 .
\end{equation}

设输出轴力矩为 $\tau_{\mathrm{out}}$、转子力矩为 $\tau$，忽略齿轮传动损失（效率 $\eta \approx 1$），
由功率守恒或杠杆关系有 $\tau_{\mathrm{out}} = i\,\tau$，即
\begin{equation}\label{eq:kt}
  \KT = \frac{K_T^{\mathrm{out}}}{i}
      = 0.3 \times \frac{187}{3591}
      \approx 0.01562\ \mathrm{N\cdot m/A}.
\end{equation}
于是转子侧力矩与力矩电流的关系为
\begin{equation}\label{eq:tau-itau}
  \tau = \KT\,\Itau = \KT \cdot \frac{20}{16384}\,\Icmd .
\end{equation}

\begin{remark}
这里把齿轮效率取为 $1$ 并不是因为它真的等于 $1$（实际上行星减速箱效率约在 $0.9$ 上下），
而是因为该效率带来的偏差最终会被第~\ref{sec:model} 节中拟合出来的系数整体吸收。
在做这类“经验模型 + 数据拟合”的工作时，把无法精确知道的物理量交给拟合参数处理，
比强行套用一个不准的常数更划算。
\end{remark}

\subsection{完整的量纲链路}

把上面几步串起来，就得到了从“发送的整型数”到“力矩”的完整映射：
\begin{equation}\label{eq:chain}
  \underbrace{\Icmd}_{\text{整型控制量}}
  \;\xrightarrow{\;\times \frac{20}{16384}\;}\;
  \underbrace{\Itau}_{\text{力矩电流 [A]}}
  \;\xrightarrow{\;\times \KT\;}\;
  \underbrace{\tau}_{\text{转子力矩 [N}\cdot\text{m]}} .
\end{equation}
\fref{eq:chain} 每一步都是一次线性标度，因此也可以合并成一个常数：
\begin{equation}
  \tau = \KT \cdot \frac{20}{16384}\,\Icmd = 1.9069\times 10^{-5}\,\Icmd \quad [\mathrm{N\cdot m}].
\end{equation}
即发送值每变化 $1$，转子力矩变化约 $1.91\times 10^{-5}\ \mathrm{N\cdot m}$；
满量程 $16384$ 对应约 $0.3124\ \mathrm{N\cdot m}$ 的转子力矩（乘回减速比即 $6.0\ \mathrm{N\cdot m}$ 输出力矩）。

\begin{table}[htbp]
  \centering
  \small
  \caption{C620 + M3508 链路中的常数}
  \label{tab:constants}
  \begin{tabular}{@{}llll@{}}
    \toprule
    量 & 符号 & 数值 & 来源 \\
    \midrule
    满量程力矩电流      & $I_{\tau,\max}$ & $20\ \mathrm{A}$           & C620 手册 \\
    控制量满量程        & $\Icmd^{\max}$  & $16384$                    & C620 手册 \\
    输出轴转矩常数      & $K_T^{\mathrm{out}}$ & $0.3\ \mathrm{N\cdot m/A}$ & M3508 手册 \\
    减速比              & $i$             & $3591/187 \approx 19.203$  & M3508 P19 \\
    转子侧力矩常数      & $\KT$           & $0.015622\ \mathrm{N\cdot m/A}$ & 式~\eqref{eq:kt} 折算 \\
    力矩电流系数        & $\CT$           & $1.99688994\times 10^{-6}$ & 式~\eqref{eq:ct} 折算 \\
    \bottomrule
  \end{tabular}
\end{table}

% ============================================================================
\section{电机功率模型}\label{sec:model}

\subsection{输入功率的构成}

无刷电机的输入功率可以按损耗机制拆成四项：
\begin{equation}\label{eq:composition}
  \Pin = \Pmech + P_{\mathrm{cu}} + P_{\mathrm{fe}} + P_{\mathrm{ctrl}},
\end{equation}
\begin{itemize}
  \item $\Pmech = \tau\Omega$：轴端输出的机械功率（$\Omega$ 为机械角速度）；
  \item $P_{\mathrm{cu}}$：绕组铜损，由力矩电流在相电阻上产生；
  \item $P_{\mathrm{fe}}$：铁损（磁滞损耗与涡流损耗），以及轴承摩擦、齿轮搅油、风阻等机械阻尼损耗；
  \item $P_{\mathrm{ctrl}}$：控制器静态功耗（MCU、栅极驱动、传感器、电调待机）。
\end{itemize}
在典型工况下机械功率占绝对主导，但其余三项正是决定“功率上限能否被精确命中”的部分，
因此必须建模而不能省略。

\subsection{机械功率与常数 9.55}

C620 回传的转速 $\omega$ 以 rpm（转/分）为单位，而功率公式要求角速度单位为 rad/s：
\begin{equation}
  \Omega = \frac{2\pi}{60}\,\omega = \frac{\omega}{9.5493}.
\end{equation}
于是
\begin{equation}\label{eq:pm}
  \Pmech = \tau\Omega = \frac{\tau\,\omega}{9.55} \qquad [\mathrm{W}],
\end{equation}
其中 $9.55 = 60/(2\pi)$ 是 rpm 到 rad/s 的换算系数。
需要注意，\fref{eq:pm} 是\emph{精确}的物理关系而\emph{不是}经验公式，
其系数 $1/9.55$ 既不需要也不可能通过数据拟合确定。这一点在第~\ref{sec:identifiability} 节会再次用到。

\subsection{损耗项：为什么是转速平方与电流平方}\label{sec:loss}

广工方案给出的经验模型形如 $\Pin = \Pmech + k_1\omega^2 + k_2\tau^2 + a$（见文献~\cite{gdutpowerformula}；需要说明的是，该公式的一手公开出处已难以查证，本文所引为其转述），
原作者也说明“查看很多资料并没有找到该式子的来源，推测为经验公式”。
事实上，这两个二次项都有清晰的物理出处，理解它们有助于判断模型何时会失效。

\paragraph{(a) 电流平方项来自铜损。}
在 $i_d = 0$ 的磁场定向控制下，电磁转矩与 $q$ 轴电流成正比：
\begin{equation}
  T_e = \tfrac{3}{2}\,p\,\psi_f\,i_q \;\equiv\; \KT\,i_q ,
\end{equation}
其中 $p$ 为极对数，$\psi_f$ 为永磁磁链（峰值约定）。三相绕组的总铜损为
\begin{equation}
  P_{\mathrm{cu}} = 3\left(\frac{i_q}{\sqrt{2}}\right)^2 R_s
  = \frac{3}{2}\,i_q^2 R_s
  = \frac{3}{2}\,\frac{R_s}{\KT^2}\,\tau^2 ,
\end{equation}
其中 $R_s$ 为单相绕组电阻。把 $\tau = \KT\Itau$ 与 \fref{eq:itau} 代入，就得到以\emph{发送值}为自变量的形式：
\begin{equation}\label{eq:k1phys}
  P_{\mathrm{cu}} = \underbrace{\frac{3}{2}\,R_s\left(\frac{20}{16384}\right)^{2}}_{k_1}\;\Icmd^2 .
\end{equation}
\fref{eq:k1phys} 说明 $k_1$ 不只是“一个拟合出来的数”，它可以反过来解读。
把部署参数 $k_1 = 1.453\times10^{-7}$ 代回，可反解出等效相电阻
\begin{equation}\label{eq:reff}
  R_{\mathrm{eff}} = \frac{k_1}{\tfrac{3}{2}\left(\frac{20}{16384}\right)^{2}}
  = \frac{1.453\times10^{-7}}{1.5\times(1.2207\times10^{-3})^{2}}
  \approx 0.065\ \Omega ,
\end{equation}
与 M3508 这类小功率无刷电机的绕组电阻量级相符。这是一个非常好用的模型自检手段：
若拟合出的 $k_1$ 反解出的电阻高达几欧姆，那一定是数据或模型出了问题，
而不是“电机比较特殊”。同时注意 $R_{\mathrm{eff}}$ 里还含着逆变桥 MOS 管的导通电阻，
所以它比手册上的纯绕组电阻略大是正常的。

\paragraph{(b) 转速平方项来自铁损与机械阻尼。}
空载（$\tau = 0$）时电机依旧要消耗功率，且该功率随转速迅速上升，来源包括：
\begin{itemize}
  \item 磁滞损耗，与频率和磁密有关，近似 $\propto f B^{n}$（$n$ 约 $1.6\sim2.2$）；
  \item 涡流损耗，$\propto f^2 B^2$；
  \item 轴承与齿轮箱的粘滞摩擦，摩擦力矩近似 $b\,\omega$，对应损耗功率 $\propto \omega^2$；
  \item 风阻与搅油损耗 $\propto \omega^3$，仅在高转速下才显著。
\end{itemize}
在 M3508 的工作转速区间内，把上述各项合并近似成一条 $\omega^2$ 抛物线是合理且方便的：
\begin{equation}\label{eq:k2phys}
  P_{\mathrm{fe}} + P_{\mathrm{damp}} \approx k_2\,\omega^2 .
\end{equation}
这也解释了开源报告中提到的一个现象——不同电机的 $k_2$ 参数不同。
原因是不同电机齿轮箱的阻尼系数不同，而阻尼力矩恰好贡献 $\omega^2$ 项。
换句话说，$k_2$ 是\emph{个体参数}，而不是整个动力套件的通用常数。

\paragraph{(c) 常数项来自控制器静态功耗。}
即使电机堵转且不输出力矩，MCU、栅极驱动、编码器、CAN 收发器与电调内部基准仍在耗电。
这部分与 $\omega$、$\Icmd$ 都无关，因此表现为常数偏置 $a$。
实测拟合得到 $a$ 在 $2\sim4\ \mathrm{W}$ 量级，与上述器件的静态功耗相符。

\subsection{模型的最终形式}

约定全文以“发送值”为自变量（与部署代码保持一致），得到
\begin{equation}\label{eq:model}
  \boxed{\;\Pin \;=\; k_2\,\omega^2 \;+\; \CT\,\Icmd\,\omega \;+\; k_1\,\Icmd^2 \;+\; a\;}
\end{equation}
其中 $\CT\Icmd\omega$ 就是 \fref{eq:pm} 的机械功率项，其系数由物理常数唯一确定：
\begin{equation}\label{eq:ct}
  \CT = \frac{\KT}{9.55}\times\frac{20}{16384}
      = \frac{0.015622}{9.55}\times1.2207\times10^{-3}
      = 1.99688994\times10^{-6}\ \ [\mathrm{W/rpm}].
\end{equation}

模型的用法如图~\ref{fig:predict} 所示：把\emph{准备下发的指令}与\emph{当前转速}
送进模型，四项分别算出再求和，就得到电机在这一刻的输入功率估计值 $\hat\Pin$。
整个过程只有几次乘加与一次平方，没有任何迭代，在 $1\ \mathrm{kHz}$ 的控制周期里开销可以忽略。

\begin{figure}[htbp]
  \centering
  \begin{tikzpicture}[>=Stealth]
    \node[fio] (in) at (0,-3.0) {$\Icmd$ 力矩电流指令\\$\omega$ 转子转速};

    \draw[thick, eblue] (2.6,-0.4) -- (2.6,-5.9);
    \draw[farr] (in.east) -- (2.6,-3.0);

    \node[fblock] (t1) at (5.3,-0.4) {$\CT\,\Icmd\,\omega$\quad\scriptsize 机械功率};
    \node[fblock] (t2) at (5.3,-2.2) {$k_2\,\omega^2$\quad\scriptsize 铁损 $+$ 阻尼};
    \node[fblock] (t3) at (5.3,-4.1) {$k_1\,\Icmd^2$\quad\scriptsize 铜损};
    \node[fblock] (t4) at (5.3,-5.9) {$a$\quad\scriptsize 控制器静态功耗};

    \draw[farr] (2.6,-0.4) -- (t1.west);
    \draw[farr] (2.6,-2.2) -- (t2.west);
    \draw[farr] (2.6,-4.1) -- (t3.west);
    \draw[farr] (2.6,-5.9) -- (t4.west);

    \node[fsum] (sum) at (9.1,-3.1) {$\Sigma$};
    \node[fblock] (out) at (11.8,-3.1) {$\hat\Pin$\\预估功率};

    \foreach \n in {t1,t2,t3,t4} { \draw[farr] (\n.east) -- (sum); }
    \draw[farr] (sum) -- (out);

    \node[flab, anchor=north west] at (-0.4,-6.8)
      {系数 $k_1,k_2,a,\CT$ 由第~\ref{sec:lsq} 节标定得到；四项之和即为预估值。};
  \end{tikzpicture}
  \caption{功率预测的实现框图。输入是“准备下发的指令”与“当前转速”，
  输出是“这样下发会消耗多少功率”。全部由四则运算构成，可以直接放进每个控制周期里算。}
  \label{fig:predict}
\end{figure}

\begin{table}[htbp]
  \centering
  \small
  \caption{功率模型参数及其记号对照}
  \label{tab:params}
  \begin{tabular}{@{}llll@{}}
    \toprule
    本文记号 & 物理含义 & 部署代码中的变量名 & 数值（本文算例）\\
    \midrule
    $k_2$   & $\omega^2$ 的系数（铁损 $+$ 阻尼） & \code{k2}                 & $1.453\times10^{-7}$ \\
    $k_1$   & $\Icmd^2$ 的系数（铜损）          & \code{a}                  & $1.230\times10^{-7}$ \\
    $a$     & 控制器静态功耗                    & \code{constant}           & $4.081$ \\
    $\CT$   & 力矩电流系数                      & \code{toque\_coefficient} & $1.99688994\times10^{-6}$ \\
    \bottomrule
  \end{tabular}
\end{table}

\begin{remark}[记号警示]
这里有两处极易踩坑的记号问题，阅读任何一份 RM 功率控制资料时都值得先确认一遍。
\begin{enumerate}
  \item 广工资料与开源报告的式 (5) 写作 $\Pin = \Pmech + k_1\omega^2 + k_2\tau^2 + a$，
  但同一份报告中后续的式 (6)(7) 与部署代码却把 $\omega^2$ 的系数记作 $k_2$、
  把电流平方项的系数记作 $k_1$——两处下标是互换的。
  本文采用与代码及式 (6)(7) 一致的那一套：$k_2$ 配 $\omega^2$，$k_1$ 配 $\Icmd^2$。
  \item 代码中被注释为 \code{// k1} 的变量 \code{a} 实际是 $\Icmd^2$ 的系数，
  被注释为 \code{// k2} 的变量实际是 $\omega^2$ 的系数。
  也就是说，代码里的 \code{a} 与公式里的常数项 $a$ \emph{不是}同一个东西。
  对照关系见表~\ref{tab:params}。
\end{enumerate}
\end{remark}

\subsection{为什么机械功率项的系数不用拟合}\label{sec:identifiability}

一个自然的问题是：既然 $k_1,k_2,a$ 都能拟合出来，为什么不把机械功率的系数也交给数据？

答案是\textbf{不可辨识}。机械功率项 $\CT\Icmd\omega$ 已由物理常数完全确定；
如果额外引入一个自由系数 $c$ 去乘 $\Icmd\omega$，那么 $c$ 与 $\CT$ 在数学上是同一个量，
数据无法区分“$\CT$ 正确、$c=1$”和“$\CT$ 偏小 $10\%$、$c=1.1$”。
更麻烦的是，$\Icmd\omega$ 与 $\Icmd^2$、$\omega^2$ 在同一批数据上往往相关性很强，
强行增加自由参数会让设计矩阵趋于病态，把本该干净的物理项误差“漏”进经验项里。

因此正确做法是：\textbf{把能精确知道的部分固定死，只让数据决定不知道的部分}。
本文模型相对早期方案最大的改进，正是把 $\tau\omega/9.55$ 按物理精确计算，
而不是用一堆二次项去“凑”出机械功率。

\subsection{模型的适用边界}

任何经验模型都有其不成立之处。以下边界在工程上必须心里有数，
其中前两条会直接导致第~\ref{sec:inverse} 节的反解失效。

\begin{enumerate}
  \item \textbf{准静态假设。} 模型描述的是稳态或慢变工况下的平均功率，
  忽略了绕组电感动态（$L\,\dd i/\dd t$）与电流环带宽限制。
  当力矩指令突变时，模型给出的瞬时功率会有偏差。
  好在控制周期（$1\ \mathrm{ms}$ 量级）远大于电流环时间常数，
  该误差在积分意义下很小。
  \item \textbf{母线电压恒定假设。} 模型没有显式包含母线电压 $U_{\mathrm{bus}}$。
  实际上同样的力矩电流，在母线电压下降时可达到的转速与功率都会变化，
  电池放电到后段时尤为明显。这一点在第~\ref{sec:future} 节作为改进方向再讨论。
  \item \textbf{温度。} 铜的电阻温度系数约为 $+0.39\%/\mathrm{K}$。
  绕组从 $25\,^{\circ}\mathrm{C}$ 升到 $100\,^{\circ}\mathrm{C}$ 时 $R_s$ 增加约 $30\%$，
  即 $k_1$ 会随温度漂移。这正是外层缓冲能量闭环必须存在的原因之一。
  \item \textbf{参数个体差异。} 如前所述，$k_2$ 与齿轮箱阻尼相关，不同电机之间不可互换；
  $k_1$ 也会因绕组批次而异。最佳实践是逐台标定，
  或至少保证用于标定的电机与装车电机同批、同保养状态。
  \item \textbf{未建模项。} 若数据中出现了无法用 $\omega^2$、$\Icmd\omega$、$\Icmd^2$ 解释的系统性偏差，
  说明模型缺少某项物理机制。此时应当去查物理，而不是继续加更高次的项。
\end{enumerate}

% ============================================================================
\section{参数辨识：最小二乘法的完整推导}\label{sec:lsq}

本节把开源资料中一笔带过的“使用 Matlab 的 \code{fit} 与 \code{fittype} 函数进行拟合”
展开成完整的数学过程。即使你今后只调用现成函数，理解这一节也能让你知道
“拟合什么时候会骗人”。整条流程可以用一张图概括，如图~\ref{fig:lsqflow} 所示；
后面的各小节就是把它逐步展开。

\begin{figure}[htbp]
  \centering
  \begin{tikzpicture}[>=Stealth, node distance=4mm]
    % 第一列：数据
    \node[fio] (d1) at (0,0)    {发送正弦力矩指令\\施加随机负载};
    \node[fio] (d2) at (0,-2.1) {采样母线电压与电流\\记录 $\omega$、$\Icmd$};
    \node[fio] (d3) at (0,-4.3) {减去已知的机械功率\\得到损耗功率 $y_i$};

    % 第二列：构造矩阵
    \node[fsoft] (x1) at (4.6,-2.1) {构造设计矩阵 $X$\\(每行 $\omega_i^2,\ {\Icmd}_i^2,\ 1$)};

    % 第三列：求解
    \node[fblock] (x2) at (9.2,-2.1) {求 $X^{\mathsf T}X$ 与 $X^{\mathsf T}y$\\解 $3\times3$ 正规方程};

    % 第四列：检验
    \node[fblock] (x3) at (9.2,-4.9) {交叉验证\\看验证集误差};
    \node[fblock] (x4) at (4.6,-6.8)  {物理自检\\$k_1$ 反解出的电阻是否合理};
    \node[fsys]  (x5) at (0,-6.8)    {装车使用\\（逐台标定）};

    \draw[farr] (d1) -- (d2);
    \draw[farr] (d2) -- (d3);
    \draw[farr] (d3.east) -- ++(0.6,0) |- (x1.west);
    \draw[farr] (x1) -- (x2);
    \draw[farr] (x2) -- (x3);
    \draw[farr] (x3.west) -- ++(-1.2,0) |- (x4.east);
    \draw[farr] (x4) -- (x5);
    \draw[fdarr] (x4.north) -- ++(0,1.35) -| (x1.south);
    \node[flab, anchor=west, text=red!65!black] at (5.0,-5.55) {不合格则重新采集};
  \end{tikzpicture}
  \caption{参数辨识的整体流程。注意最后一道“物理自检”：
  由于 $k_1$ 可以反解出等效相电阻，拟合结果是否可信有一个与数据无关的独立判据。}
  \label{fig:lsqflow}
\end{figure}

\subsection{把辨识问题写成线性回归}

\fref{eq:model} 中待定的参数只有三个：$k_2,k_1,a$。
把它们记成参数向量 $\theta = (k_2,\,k_1,\,a)^{\mathsf T}$。模型对 $\theta$ 是\emph{线性}的，
这是一个非常好的性质：线性模型的最小二乘问题有闭式解，不需要迭代（标准参考见文献~\cite{golub2013matrix,bjorck1996numerical}）。

先把已知的机械功率项移到左边。定义
\begin{equation}
  y_i \;\triangleq\; \Pin^{(i)} - \CT\,\Icmd^{(i)}\,\omega^{(i)},
\end{equation}
即“实测输入功率减去已知的机械功率”，其物理含义就是\emph{总损耗功率}。
于是对第 $i$ 组样本（$i=1,\dots,N$）有
\begin{equation}\label{eq:regmodel}
  y_i = k_2\bigl(\omega^{(i)}\bigr)^2 + k_1\bigl(\Icmd^{(i)}\bigr)^2 + a + \varepsilon_i ,
\end{equation}
其中 $\varepsilon_i$ 表示模型误差与测量噪声之和。写成矩阵形式：
\begin{equation}\label{eq:designmat}
  X = \begin{bmatrix}
    (\omega^{(1)})^2 & (\Icmd^{(1)})^2 & 1\\[2pt]
    (\omega^{(2)})^2 & (\Icmd^{(2)})^2 & 1\\[2pt]
    \vdots & \vdots & \vdots\\[2pt]
    (\omega^{(N)})^2 & (\Icmd^{(N)})^2 & 1
  \end{bmatrix}\in\mathbb{R}^{N\times 3},
  \quad
  y=\begin{bmatrix} y_1\\ \vdots\\ y_N\end{bmatrix},
  \quad
  \varepsilon=\begin{bmatrix} \varepsilon_1\\ \vdots\\ \varepsilon_N\end{bmatrix},
\end{equation}
则 \fref{eq:regmodel} 整体就是 $y = X\theta + \varepsilon$。
矩阵 $X$ 的每一列称为一个\emph{回归量}（regressor）。

\begin{table}[htbp]
  \centering
  \small
  \caption{辨识问题中各个矩阵的形状与含义}
  \label{tab:matrixshape}
  \begin{tabular}{@{}llll@{}}
    \toprule
    符号 & 形状 & 含义 & 本问题的取值 \\
    \midrule
    $X$            & $N\times 3$ & 设计矩阵（每行一个样本） & $N$ 为样本数，列数为 $3$ \\
    $\theta$       & $3\times 1$ & 待估参数 $(k_2,\;k_1,\;a)^{\mathsf T}$ & $3$ \\
    $y$            & $N\times 1$ & 每组的损耗功率 & $N$ \\
    $\varepsilon$  & $N\times 1$ & 残差 & $N$ \\
    $X^{\mathsf T}X$ & $3\times 3$ & 回归量的内积矩阵 & 与 $N$ 无关 \\
    $X^{\mathsf T}y$ & $3\times 1$ & 回归量与观测量的内积 & 与 $N$ 无关 \\
    \bottomrule
  \end{tabular}
\end{table}

\begin{figure}[htbp]
  \centering
  \begin{tikzpicture}[>=Stealth, x=1cm, y=1cm]
    % X^T
    \draw[fill=eblue!12, draw=eblue, thick] (0,0.99) rectangle (2.6,1.61);
    \node[font=\small] at (1.3,1.3) {$X^{\mathsf T}$};
    \node[flab, anchor=north] at (1.3,0.91) {$3\times N$};

    \node[font=\large] at (3.1,1.3) {$\times$};

    % X
    \draw[fill=codegreen!12, draw=codegreen, thick] (3.6,0) rectangle (4.22,2.6);
    \node[font=\small, rotate=90] at (3.91,1.3) {$X$};
    \node[flab, anchor=north] at (3.91,-0.08) {$N\times3$};

    \node[font=\large] at (4.7,1.3) {$=$};

    % X^T X
    \draw[fill=eblue!12, draw=eblue, thick] (5.2,0.6) rectangle (6.6,2.0);
    \node[font=\small] at (5.9,1.3) {$X^{\mathsf T}X$};
    \node[flab, anchor=north] at (5.9,0.52) {$3\times3$};

    \node[font=\large] at (7.1,1.3) {$\times$};

    % theta
    \draw[fill=red!10, draw=red!65!black, thick] (7.6,0.6) rectangle (8.22,2.0);
    \node[font=\small, rotate=90] at (7.91,1.3) {$\theta$};
    \node[flab, anchor=north] at (7.91,0.52) {$3\times1$};

    \node[font=\large] at (8.7,1.3) {$=$};

    % X^T y
    \draw[fill=red!10, draw=red!65!black, thick] (9.2,0.6) rectangle (9.82,2.0);
    \node[font=\small, rotate=90] at (9.51,1.3) {$X^{\mathsf T}y$};
    \node[flab, anchor=north] at (9.51,0.52) {$3\times1$};

    % 注释
    \node[flab, anchor=north west, align=left] at (0,-0.55)
      {$N$ 可达上万，但两个内积矩阵只与列的个数（$3$）有关，};
    \node[flab, anchor=north west, align=left] at (0,-1.05)
      {因此最终要解的始终只是一个 $3\times3$ 的线性方程组。};
  \end{tikzpicture}
  \caption{正规方程的矩阵形状。注意“大矩阵”$X$ 在被利用完之后就消失了：
  真正需要求解的只是 $X^{\mathsf T}X\;\theta=X^{\mathsf T}y$ 这个 $3\times3$ 的小方程组。}
  \label{fig:matrixshape}
\end{figure}

\begin{remark}[为什么要先减掉机械功率]
这一步常被忽略，但很关键。若直接对总功率 $\Pin$ 拟合，
目标函数中占绝对主导的是 $\CT\Icmd\omega$（典型工况下可达几十瓦），
而三个待估参数对应的损耗只有几瓦。优化器会把绝大部分“注意力”花在拟合机械功率上，
损耗参数的相对精度会差得离谱。先减去已知的物理项，
等价于把待拟合量的动态范围压缩了一个数量级，
是这类“灰箱辨识”中的标准技巧。
\end{remark}

\subsection{最小二乘准则}

最小二乘的思想很朴素：选一组参数，使所有样本的\emph{误差平方和}最小。定义准则函数
\begin{equation}\label{eq:cost}
  S(\theta) \;=\; \sum_{i=1}^{N}\varepsilon_i^2
  \;=\;\sum_{i=1}^{N}\Bigl(y_i - k_2\bigl(\omega^{(i)}\bigr)^2 - k_1\bigl(\Icmd^{(i)}\bigr)^2 - a\Bigr)^{2}.
\end{equation}
为什么用平方而不是绝对值？有四个理由：
\begin{enumerate}
  \item 平方处处可导，可以用微积分求极值，从而得到闭式解；
  \item $S(\theta)$ 关于 $\theta$ 是凸二次函数，\emph{局部极小就是全局极小}，
  不存在陷入局部最优的风险；
  \item 平方对大的偏差惩罚更重；
  \item 当 $\varepsilon_i$ 独立同分布且服从高斯分布时，最小化 $S$ 等价于极大似然估计，
  即在这种噪声假设下最小二乘是最优估计。
\end{enumerate}

\subsection{用求导求极值：一阶必要条件}\label{sec:firstorder}

下面把求导过程一步步写清楚，全程只用三条基本法则：
\begin{align}
  \frac{\partial}{\partial x}\bigl[u(x)+v(x)\bigr] &= \frac{\partial u}{\partial x}+\frac{\partial v}{\partial x},\label{eq:rule1}\\
  \frac{\partial}{\partial x}\bigl[c\,u(x)\bigr] &= c\,\frac{\partial u}{\partial x}\quad(c\ \text{为常数}),\label{eq:rule2}\\
  \frac{\partial}{\partial x}\bigl[u(x)^2\bigr] &= 2u\,\frac{\partial u}{\partial x}.\label{eq:rule3}
\end{align}

把 $S$ 中的每一项看成待求导的复合函数。记第 $i$ 项的残差为
\begin{equation}
  r_i(\theta) \;\triangleq\; y_i - k_2\bigl(\omega^{(i)}\bigr)^2 - k_1\bigl(\Icmd^{(i)}\bigr)^2 - a ,
\end{equation}
则 $S=\sum_i r_i^2$。对任一参数 $\theta_j\in\{k_2,k_1,a\}$，由 \fref{eq:rule1}--\fref{eq:rule3} 得
\begin{equation}
  \frac{\partial S}{\partial \theta_j}
  = \sum_{i=1}^{N}\frac{\partial}{\partial \theta_j}\bigl(r_i^2\bigr)
  = \sum_{i=1}^{N} 2\,r_i\,\frac{\partial r_i}{\partial \theta_j},
\end{equation}
而 $r_i$ 对三个参数的偏导数分别是
\begin{equation}
  \frac{\partial r_i}{\partial k_2} = -\bigl(\omega^{(i)}\bigr)^2,\qquad
  \frac{\partial r_i}{\partial k_1} = -\bigl(\Icmd^{(i)}\bigr)^2,\qquad
  \frac{\partial r_i}{\partial a} = -1 .
\end{equation}
在极小值处三个偏导数必须同时为零（多元函数取极值的必要条件）。
把公共常数 $-2$ 约去，得到三个条件：
\begin{equation}\label{eq:orth}
  \boxed{\;
  \sum_{i=1}^{N} r_i\bigl(\omega^{(i)}\bigr)^2 = 0,\qquad
  \sum_{i=1}^{N} r_i\bigl(\Icmd^{(i)}\bigr)^2 = 0,\qquad
  \sum_{i=1}^{N} r_i = 0 .
  \;}
\end{equation}

\fref{eq:orth} 有一个漂亮的几何意义：\textbf{最优残差向量与每一个回归量都正交}。
也就是说，最小二乘解把数据中“能被这三个回归量解释的部分”全部榨干了，
剩下的残差与它们毫无关系。这个“正交性”正是最小二乘理论的骨架。

把 $r_i$ 的定义代回 \fref{eq:orth} 并展开，就得到三个关于 $(k_2,k_1,a)$ 的\emph{线性}方程。
为书写简洁，下面把 $(\cdot)^{(i)}$ 简写为 $(\cdot)_i$：
\begin{align}
  k_2\sum_i \omega_i^4 \;+\; k_1\sum_i \omega_i^2{\Icmd}_i^2 \;+\; a\sum_i \omega_i^2
    &= \sum_i \omega_i^2\,y_i,\label{eq:ne1}\\
  k_2\sum_i \omega_i^2{\Icmd}_i^2 \;+\; k_1\sum_i {\Icmd}_i^4 \;+\; a\sum_i {\Icmd}_i^2
    &= \sum_i {\Icmd}_i^2\,y_i,\label{eq:ne2}\\
  k_2\sum_i \omega_i^2 \;+\; k_1\sum_i {\Icmd}_i^2 \;+\; N\,a
    &= \sum_i y_i .\label{eq:ne3}
\end{align}
\fref{eq:ne1}--\fref{eq:ne3} 称为\textbf{正规方程}（normal equations）。
注意它的结构：左边是“回归量之间的内积”乘参数，右边是“回归量与观测量之间的内积”。
只要把数据代入算出这 6 个求和量，就得到一个三元一次方程组，
用消元法或克拉默法则即可解出 $\hat k_2,\hat k_1,\hat a$。

\subsection{矩阵形式：几行推导代替一页求和}

用矩阵记号可以把上面的过程压缩到三行，并且顺便得到更一般的结论。
把 \fref{eq:cost} 写成向量范数的形式：
\begin{equation}\label{eq:costmat}
  S(\theta) = \lVert y - X\theta\rVert^2 = (y-X\theta)^{\mathsf T}(y-X\theta).
\end{equation}
展开括号。注意 $\theta^{\mathsf T}X^{\mathsf T}y$ 与 $y^{\mathsf T}X\theta$ 都是标量且互为转置，
因此两者相等，可以合并为一项：
\begin{equation}\label{eq:expand}
  S(\theta) = y^{\mathsf T}y - 2\,\theta^{\mathsf T}X^{\mathsf T}y + \theta^{\mathsf T}X^{\mathsf T}X\,\theta .
\end{equation}
标量对向量求导有两条常用公式：
\begin{equation}\label{eq:matrixderiv}
  \frac{\partial\,(a^{\mathsf T}\theta)}{\partial\theta} = a,
  \qquad
  \frac{\partial\,(\theta^{\mathsf T}A\theta)}{\partial\theta} = (A+A^{\mathsf T})\,\theta
  \;\overset{A=A^{\mathsf T}}{=}\; 2A\theta .
\end{equation}
把它用到 \fref{eq:expand} 的三项上：第一项 $y^{\mathsf T}y$ 与 $\theta$ 无关，导数为零；
第二项给出 $-2X^{\mathsf T}y$；第三项给出 $2X^{\mathsf T}X\theta$。于是
\begin{equation}\label{eq:gradS}
  \frac{\partial S}{\partial\theta}
  = -2X^{\mathsf T}y + 2X^{\mathsf T}X\,\theta .
\end{equation}
令 \fref{eq:gradS} 为零，立即得到
\begin{equation}\label{eq:normal}
  \boxed{\;X^{\mathsf T}X\,\hat\theta = X^{\mathsf T}y\;}
  \qquad\Longrightarrow\qquad
  \boxed{\;\hat\theta = (X^{\mathsf T}X)^{-1}X^{\mathsf T}y\;}
\end{equation}
这与 \fref{eq:ne1}--\fref{eq:ne3} 完全等价，只是写法更紧凑、更容易推广到更多参数的情形。

把 \fref{eq:normal} 的两个矩阵完整写出来，就是下面这个样子。
看到它就能明白为什么我说“最后只需要解一个小方程组”：
矩阵里的每一个元素都是一个求和量，$N$ 再大，矩阵也只有 $3\times3$。
\begin{equation}\label{eq:normalexpanded}
  \underbrace{
  \begin{bmatrix}
    \sum_i \omega_i^{4} & \sum_i \omega_i^{2}{\Icmd}_i^{2} & \sum_i \omega_i^{2}\\[6pt]
    \sum_i \omega_i^{2}{\Icmd}_i^{2} & \sum_i {\Icmd}_i^{4} & \sum_i {\Icmd}_i^{2}\\[6pt]
    \sum_i \omega_i^{2} & \sum_i {\Icmd}_i^{2} & N
  \end{bmatrix}}_{X^{\mathsf T}X\ (3\times3)}
  \underbrace{
  \begin{bmatrix} k_2\\[4pt] k_1\\[4pt] a\end{bmatrix}}_{\hat\theta\ (3\times1)}
  =
  \underbrace{
  \begin{bmatrix} \sum_i \omega_i^{2}y_i\\[6pt] \sum_i {\Icmd}_i^{2}y_i\\[6pt] \sum_i y_i\end{bmatrix}
  }_{X^{\mathsf T}y\ (3\times1)}
\end{equation}

注意 $X^{\mathsf T}X$ 是\emph{对称}的，三个对角元分别是
$\sum\omega^4$、$\sum\Icmd^4$、$N$，
也就是说它只由数据的四个求和量（$\sum\omega^4$、$\sum\Icmd^4$、
$\sum\omega^2\Icmd^2$、$\sum\omega^2$、$\sum\Icmd^2$）决定。
实际编程时不需要真的构造 $X$，只需遍历一次数据把这几个量累加起来，
再解一个三元一次方程组即可——这正是它适合放进单片机的根本原因。

顺带说明为什么这个驻点一定是极小值：对 \fref{eq:gradS} 再求一次导，得到海森矩阵
\begin{equation}
  \frac{\partial^2 S}{\partial\theta\,\partial\theta^{\mathsf T}} = 2X^{\mathsf T}X \succeq 0 ,
\end{equation}
它是半正定的（因为二次型
$\theta^{\mathsf T}X^{\mathsf T}X\theta=\lVert X\theta\rVert^2\ge0$ 恒非负），
所以 $S(\theta)$ 是凸函数，驻点即全局最小点。

\subsection{解是否存在、是否唯一}

\fref{eq:normal} 中的求逆只在 $X^{\mathsf T}X$ 可逆时成立。而
\begin{equation}
  X^{\mathsf T}X\ \text{可逆} \iff X\ \text{列满秩} \iff
  \text{三个回归量在这批数据上线性无关}.
\end{equation}
下面是两个典型的“不可辨识”场景：
\begin{itemize}
  \item 如果所有实验都在\emph{同一个恒定转速}下完成，那么 $\omega_i^2$ 对所有 $i$ 都等于同一个常数，
  第一列就是第三列（全 $1$ 列）的倍数。此时 $k_2$ 与 $a$ 无法分开：
  把 $k_2$ 调大、$a$ 调小可以给出完全相同的预测，数据没有能力区分它们。
  \item 如果整个实验过程中力矩指令始终不变，则 ${\Icmd}_i^2$ 恒定，$k_1$ 与 $a$ 同样不可分。
\end{itemize}
因此\textbf{实验设计必须让两个回归量都充分变化}。
这正好解释了数据集里为什么既要有变转速的正弦、斜坡激励，又要有随机负载：
这不是为了把数据量做大，而是为了让 $X$ 的各列去相关，从而让参数可辨识。

即使 $X$ 严格列满秩，若列之间高度相关（近似共线），参数的方差也会被放大。
在噪声独立同分布、方差为 $\sigma^2$ 的假设下，最小二乘估计的协方差为
\begin{equation}\label{eq:cov}
  \operatorname{Cov}(\hat\theta) = \sigma^2 (X^{\mathsf T}X)^{-1}.
\end{equation}
列相关性越强，$(X^{\mathsf T}X)^{-1}$ 的对角元越大，估计越不可信。
工程上可以计算方差膨胀因子（VIF）或 $X^{\mathsf T}X$ 的条件数来做体检。

\subsection{数值条件：一个容易被忽视的坑}

$X$ 的三列在量纲上差异巨大：$\omega$ 可达 $10^4\ \mathrm{rpm}$，故 $\omega^2$ 可达 $10^8$；
$\Icmd$ 可达 $16384$，故 $\Icmd^2$ 可达 $2.7\times10^8$；而第三列恒为 $1$。
于是 $X^{\mathsf T}X$ 的元素跨越十几个数量级，其条件数是 $X$ 的条件数的平方：
\begin{equation}
  \kappa(X^{\mathsf T}X) = \kappa(X)^{2}.
\end{equation}
在双精度（相对精度约 $10^{-16}$）下这通常还能求解，但若用单精度做同样的事，
有效数字会被吃掉大半。三条实用的改进措施：

\begin{enumerate}
  \item \textbf{无量纲化。} 取参考量 $\omega_{\mathrm{ref}}$、$I_{\mathrm{ref}}$
  （例如各自的最大值），令 $\tilde\omega=\omega/\omega_{\mathrm{ref}}$、
  $\tilde I=\Icmd/I_{\mathrm{ref}}$，在新变量上拟合出 $\tilde k_2,\tilde k_1,\tilde a$，
  再按量纲换算回去：
  \begin{equation}\label{eq:rescaling}
    k_2 = \frac{\tilde k_2}{\omega_{\mathrm{ref}}^{2}},\qquad
    k_1 = \frac{\tilde k_1}{I_{\mathrm{ref}}^{2}},\qquad
    a = \tilde a .
  \end{equation}
  这样 $X^{\mathsf T}X$ 的所有元素都落在同一量级上。
  \item \textbf{中心化。} 改用 $\omega_i^2-\overline{\omega^2}$ 与
  ${\Icmd}_i^2-\overline{\Icmd^2}$ 作为回归量，使它们与截距列正交，
  $k_2,k_1$ 的估计与 $a$ 解耦。代价是此时 $a$ 的含义变为“平均工况下的损耗”，
  需要按定义换算回来。
  \item \textbf{用矩阵分解而不是显式求逆。} 数值上应当通过 QR 分解或奇异值分解求解
  \fref{eq:normal}，而不是先算 $(X^{\mathsf T}X)^{-1}$ 再乘 $X^{\mathsf T}y$。
  MATLAB 的 \code{\\} 运算符与 \code{fit} 内部正是这样做的，
  所以在 MATLAB 里通常不会遇到精度问题；但把同样的公式手写进单片机就是另一回事了。
\end{enumerate}

\subsection{当最小二乘不够用：加权、鲁棒与误差变量}

实际数据很少满足最小二乘的理想假设。以下三种修正值得知道。

\paragraph{加权最小二乘（异方差噪声）。}
若不同样本的噪声方差不同，应当给更可靠的样本更大的权重。极小化
\begin{equation}
  S_W(\theta)=\sum_{i=1}^{N} w_i\,r_i^{2}
\end{equation}
得到的解为
\begin{equation}\label{eq:wls}
  \hat\theta_W = (X^{\mathsf T}WX)^{-1}X^{\mathsf T}Wy ,
  \qquad W=\operatorname{diag}(w_1,\dots,w_N).
\end{equation}
取 $w_i = 1/\sigma_i^{2}$ 时，该估计在所有线性无偏估计中方差最小（高斯--马尔可夫定理）。
在本问题中一个很实际的做法是：把 INA226 采样接近满量程而出现削顶的样本
（开源报告图 1 的注记正说明了这一现象）以及转速接近零的样本权重取 $0$，直接剔除。

\paragraph{鲁棒回归（异常值）。}
最小二乘的“崩溃点”为零，少数几个离谱样本就能把整条曲线拽偏。
若怀疑存在异常值，可用 Huber 损失~\cite{huber2009robust} 或迭代重加权最小二乘（IRLS）：
\begin{equation}
  \theta^{(t+1)} = \bigl(X^{\mathsf T}W^{(t)}X\bigr)^{-1}X^{\mathsf T}W^{(t)}y,
  \qquad w_i^{(t)} = \frac{\psi\bigl(r_i^{(t)}\bigr)}{r_i^{(t)}},
\end{equation}
Huber 的 $\psi$ 在 $|r|\le\delta$ 时等于 $r$、超出后饱和，因此异常点的权重会被自动压小。
形式上它和加权最小二乘是同一个公式，只是权重在迭代中不断更新。

\paragraph{误差变量（回归量本身含噪声）。}
本问题中 $\Icmd$ 是我们自己发出去的数，精确已知；但 $\omega$ 来自编码器，
存在量化与测量噪声。带噪声的回归量会破坏“噪声与回归量不相关”这一前提，
使最小二乘出现\textbf{衰减偏置}：由于
$\mathbb{E}\bigl[\omega_{\mathrm{meas}}^{2}\bigr]=\omega^{2}+\sigma_\omega^{2}$，
$\omega^2$ 这一列被系统性地抬高，估计出的 $k_2$ 会偏小，
并且该偏置会顺着正规方程污染 $k_1$ 与 $a$。
可行的缓解手段包括：提高测速精度（滑动平均、或降低采样率换取分辨率）、
使用总体最小二乘／正交距离回归，或先估计出 $\sigma_\omega^{2}$ 再做偏置修正。
对本文这种“够用就行”的工程场景，保证标定时的测速条件与实车一致即可。

\subsection{拟合优度与交叉验证}

拟合完成后至少应报告两个指标：
\begin{equation}
  R^{2} = 1-\frac{\sum_i (y_i-\hat y_i)^2}{\sum_i (y_i-\bar y)^2},
  \qquad
  \mathrm{RMSE} = \sqrt{\frac{1}{N}\sum_{i=1}^{N}(y_i-\hat y_i)^{2}} .
\end{equation}
但要注意，只有三个参数的模型 $R^2$ 几乎必然“好看”，
它并不能说明模型在没见过的工况下是否成立。
更有意义的是\emph{交叉验证}：把若干组独立采集的数据
（不同激励形式、不同负载）轮流作为训练集与验证集，
观察参数是否稳定、验证集 RMSE 是否与训练集接近。
开源实现正是通过“不同组数据交叉验证仿真效果”来挑选最终参数的，
这一步比单纯看 $R^2$ 重要得多。

\subsection{拟合结果与自检}

\begin{table}[htbp]
  \centering
  \small
  \caption{两组参数辨识结果对比（$\omega$ 单位 rpm，$\Icmd$ 为无量纲控制量）}
  \label{tab:fit}
  \begin{tabular}{@{}lccc@{}}
    \toprule
    来源 & $\omega^2$ 系数 $k_2$ & $\Icmd^2$ 系数 $k_1$ & 常数项 $a$ \\
    \midrule
    \code{power\_model.m} 中的拟合示例 & $2.110\times10^{-7}$ & $9.805\times10^{-8}$ & $2.138$ \\
    部署代码 \code{chassis\_power\_control.c} & $1.453\times10^{-7}$ & $1.230\times10^{-7}$ & $4.081$ \\
    \bottomrule
  \end{tabular}
\end{table}

表~\ref{tab:fit} 中两组参数差异可见，这是完全正常的：它们来自不同的电机样本、
不同的负载条件与不同的数据组。这也再次印证了关于 $k_2$ 个体差异的讨论。
真正需要保证的是——\textbf{部署用的参数必须来自同一台（或同批同状态）的电机}。

用 \fref{eq:k1phys} 反解等效电阻可以给这两组参数做一次廉价的自检：
第一组 $R_{\mathrm{eff}}\approx0.044\ \Omega$，第二组 $R_{\mathrm{eff}}\approx0.065\ \Omega$，
都落在合理区间内。若某次拟合给出 $k_1$ 为 $10^{-5}$ 量级，
那对应的等效电阻将是几欧姆，此时应当优先怀疑数据而不是电机。

% ============================================================================
\section{给定功率反解力矩}\label{sec:inverse}

有了模型，就可以回答第~\ref{sec:convention} 节提出的核心问题：
\emph{在当前转速下，为了恰好消耗给定功率 $P^{\ast}$，应该发送多大的控制量？}
这一步是整个算法的执行环节，也是开源资料中讨论最充分、却也最容易踩坑的地方。

\subsection{反解方程的建立}

设当前转速为 $\omega$，目标功率为 $P^{\ast}$，未知量为发送值 $x=\Icmd$。
把 $x$ 代入 \fref{eq:model} 并令 $\Pin = P^{\ast}$：
\begin{equation}\label{eq:invsetup}
  k_2\omega^2 + \CT\omega\,x + k_1x^2 + a = P^{\ast}.
\end{equation}
整理成标准的一元二次方程
\begin{equation}\label{eq:invquad}
  \boxed{\;A x^2 + B x + C = 0\;}
  \qquad\text{其中}\qquad
  \begin{cases}
    A = k_1,\\[2pt]
    B = \CT\,\omega,\\[2pt]
    C = k_2\omega^2 + a - P^{\ast}.
  \end{cases}
\end{equation}
注意 $B$ 与 $\omega$ 成正比、符号相同；$A=k_1>0$ 恒成立（它是电阻，必然为正）。
$C$ 的符号则取决于目标功率与该转速下\emph{空载损耗}的比较。

为后文方便，定义两个辅助量。\textbf{空载损耗}（$x=0$ 时的输入功率）
\begin{equation}\label{eq:noload}
  P_0(\omega) \;\triangleq\; k_2\omega^2 + a ,
\end{equation}
以及\textbf{最小可达功率}（抛物线顶点处的功率）
\begin{equation}\label{eq:pmin}
  P_{\min}(\omega) \;\triangleq\; P_0(\omega) - \frac{B^2}{4A}
  = P_0(\omega) - \frac{\CT^2\omega^2}{4k_1}.
\end{equation}

\subsection{判别式与解的几何图像}

一元二次方程的判别式为
\begin{equation}\label{eq:disc}
  \Delta = B^2 - 4AC .
\end{equation}
实根存在的条件是 $\Delta\ge0$，把它翻译回物理量：
\begin{equation}\label{eq:cond}
  \Delta \ge 0
  \iff P^{\ast} \ge P_{\min}(\omega)
  \iff P^{\ast} \ge \Bigl(k_2 - \frac{\CT^2}{4k_1}\Bigr)\omega^{2} + a .
\end{equation}
把数值代进去会看到一个有意思的结果。按表~\ref{tab:params} 的参数，
\begin{equation}
  \frac{\CT^{2}}{4k_1} = \frac{(1.9969\times10^{-6})^{2}}{4\times1.23\times10^{-7}}
  \approx 8.105\times10^{-6}
  \;\gg\; k_2 = 1.453\times10^{-7},
\end{equation}
因此 \fref{eq:pmin} 中 $\omega^2$ 的系数为负，且绝对值远大于 $k_2$：
\begin{equation}\label{eq:pminnum}
  P_{\min}(\omega) \approx 4.081 - 7.96\times10^{-6}\,\omega^{2}\quad[\mathrm{W}].
\end{equation}
令 $P_{\min}(\omega)=0$ 得到 $\omega\approx716\ \mathrm{rpm}$。也就是说：

\begin{itemize}
  \item 当 $\omega>716\ \mathrm{rpm}$ 时，$P_{\min}<0$。由于我们只关心非负的功率预算，
  \textbf{判别式恒为非负，反解永远有解}；
  \item 只有\emph{低速且功率预算极低}（几百转、几瓦）时，才可能出现 $\Delta<0$。
\end{itemize}
这个结论很实用：很多实现里那个“判别式小于零就跳过”的分支，
在正常行驶工况下几乎永远不会被触发；真正的风险不在判别式，而在下面要讲的\emph{根的符号}。

\begin{figure}[htbp]
  \centering
  \begin{tikzpicture}
    \begin{axis}[
      width=0.88\textwidth, height=6.4cm,
      xlabel={力矩电流控制值 $\Icmd$},
      ylabel={电机输入功率 $\Pin$ (W)},
      xmin=-4000, xmax=12000, ymin=-60, ymax=190,
      grid=both, grid style={gray!18},
      legend style={at={(0.98,0.03)}, anchor=south east, draw=gray!40,
                    fill=white, fill opacity=0.9, text opacity=1, font=\small},
      tick label style={font=\small},
      label style={font=\small},
      axis line style={gray!60}
    ]
      \addplot[domain=-4000:12000, samples=120, thick, eblue]
        {1.23e-7*x^2 + 0.0119813*x + 9.3118};
      \addlegendentry{$\omega = 6000$ rpm}
      \addplot[domain=-4000:12000, samples=120, thick, codegreen, dashed]
        {1.23e-7*x^2 + 0.0059907*x + 5.3887};
      \addlegendentry{$\omega = 3000$ rpm}
      \addplot[gray, thick, dotted] coordinates {(-4000,40) (12000,40)};
      \addlegendentry{功率预算 $P^{\ast} = 40$ W}
      \addplot[only marks, mark=*, mark size=1.8pt, black]
        coordinates {(2497,40) (5898,40)};
      \node[anchor=south west, font=\scriptsize, text=black] at (axis cs:2550,42)
        {$x_+\!\approx\!2497$};
      \node[anchor=south east, font=\scriptsize, text=black] at (axis cs:5880,42)
        {$x_+\!\approx\!5898$};
    \end{axis}
  \end{tikzpicture}
  \caption{同一功率预算在不同转速下对应的控制量。转速越低，允许的控制量越大，
  这正是“固定力矩阈值”方案失效的根本原因：功率与控制量的换算关系随转速剧烈变化。}
  \label{fig:powercurve}
\end{figure}

\subsection{两个根都有物理意义}

当 $\Delta>0$ 时，方程有两个实根
\begin{equation}\label{eq:roots}
  x_{\pm} = \frac{-B \pm \sqrt{\Delta}}{2A}
          = \frac{-\CT\omega \pm \sqrt{\Delta}}{2k_1}.
\end{equation}
由韦达定理可以得到两个非常有用的关系：
\begin{equation}\label{eq:vieta}
  x_+ + x_- = -\frac{B}{A} = -\frac{\CT\omega}{k_1},
  \qquad
  x_+ x_- = \frac{C}{A} = \frac{k_2\omega^2 + a - P^{\ast}}{k_1}.
\end{equation}
把 \fref{eq:vieta} 与 \fref{eq:noload} 结合起来看，可以完整刻画根的情况。
设 $\omega>0$（否则整体取负号即可），则 $B>0$，因而两根之和恒为负。

\begin{description}
  \item[情形 1：$P^{\ast} > P_0(\omega)$，即目标功率高于空载损耗。]
  此时 $C<0$，两根之积为负，说明\textbf{一正一负}。
  正根对应“驱动”，负根对应“回馈制动”。这是正常行驶时的情形。
  \item[情形 2：$P^{\ast} = P_0(\omega)$。] 此时 $C=0$，一个根为 $0$，另一个为 $-B/A<0$。
  零根意味着“不发力，让电机自由空转”。
  \item[情形 3：$P_{\min}(\omega) \le P^{\ast} < P_0(\omega)$。]
  此时 $C>0$，两根之积为正，而和恒为负，说明\textbf{两根都为负}——
  只有两个制动解，不存在正向驱动解。
  \item[情形 4：$P^{\ast} < P_{\min}(\omega)$。] $\Delta<0$，无实根，目标功率不可达。
\end{description}

情形 3 的物理含义值得单独说明。空载损耗 $P_0(\omega)$ 是电机“不发力”时的固有消耗：
即使输出力矩为零，铁损、齿槽阻尼与控制器静态功耗也照样存在。
因此，若要求电机把消耗降到 $P_0(\omega)$ 以下，唯一的办法是让它\emph{轻微回馈制动}——
用一个小的负力矩把机械能送回母线。这就是为什么会出现两个负根。

这也解释了开源报告中观察到的现象：

\begin{quote}
“在某一时刻若转速为正，若让它消耗同一功率，其输出的正的加速力矩绝对值
要小于负的减速力矩的绝对值。也就是说，在减速时，相同功率，我们可以输出更高的减速力矩，
这和电动机反电动势的原理是一致的。”
\end{quote}

用 \fref{eq:vieta} 可以直接给出这个结论：两根异号且和恒为负，
因此负根的绝对值必然大于正根的绝对值，即 $|x_-|>|x_+|$。

\subsection{根的选择规则}

既然两根都是方程的合法解，就必须有一个选择准则。这里最容易犯的错误，
恰恰是开源报告作者最初踩过的坑：

\begin{quote}
\itshape
“在起初笔者直接使用正号的解作为目标力矩的大小，
然后依据原始力矩的正负符号对其重新赋值，在实际效果表现很差。”
\end{quote}

\emph{先取正根再用原始符号强行赋值}这种做法是错的，
因为它把一个本该由模型决定的信息（应该在驱动侧还是制动侧工作）交给了符号翻转，
而不是让模型去判断。正确的准则是：

\begin{quote}
\textbf{以缩放前的原始控制量 $\Icmd$ 的符号为准，选择与之同号的那一个根。}
\end{quote}

即
\begin{equation}\label{eq:rootselect}
  x' =
  \begin{cases}
    \dfrac{-B + \sqrt{\Delta}}{2A}, & \Icmd > 0 \quad(\text{原始意图为驱动}),\\[10pt]
    \dfrac{-B - \sqrt{\Delta}}{2A}, & \Icmd < 0 \quad(\text{原始意图为制动}).
  \end{cases}
\end{equation}
这样做的理由很直白：它保证功率限制\textbf{只改变力矩的大小、不改变力矩的方向}。
若不这样做，就可能把驱动指令变成制动指令，或者反过来，
在实车上直接表现为车身抖动与方向突变。

\subsection{数值算例}

下面把一次完整的反解过程手算一遍，参数取自表~\ref{tab:params}。

\begin{table}[htbp]
  \centering
  \small
  \caption{反解算例（$k_1=1.23\times10^{-7}$，$k_2=1.453\times10^{-7}$，$a=4.081$，$\CT=1.9969\times10^{-6}$）}
  \label{tab:example}
  \begin{tabular}{@{}lccc@{}}
    \toprule
    步骤 & 算例 A（常规） & 算例 B（低于空载损耗） & 算例 C（低速不可达） \\
    \midrule
    转速 $\omega$ (rpm)              & $6000$ & $1000$ & $400$ \\
    目标功率 $P^{\ast}$ (W)          & $40$   & $4$    & $2$ \\
    空载损耗 $P_0(\omega)$ (W)       & $9.312$ & $4.226$ & $4.104$ \\
    $B=\CT\omega$                    & $0.011981$ & $0.001997$ & $0.000799$ \\
    $C=k_2\omega^2+a-P^{\ast}$       & $-30.688$ & $0.2263$ & $2.1042$ \\
    $\Delta=B^2-4AC$                 & $1.587\times10^{-4}$ & $3.876\times10^{-6}$ & $-3.973\times10^{-7}$ \\
    判别结论                          & 一正一负 & 两根皆负 & 无实根 \\
    $x_+$                            & $2497$ & $-114.1$ & --- \\
    $x_-$                            & $-9.99\times10^{4}$ & $-1.612\times10^{4}$ & --- \\
    原始控制量 $\Icmd$ 符号          & $+$ & $+$ & $+$ \\
    按 \fref{eq:rootselect} 取值     & $x_+=2497$ & $x_+=-114$（\textcolor{red}{异常}） & --- \\
    \midrule
    回代校验 $\Pin(x')$ (W)          & $39.99$ & $4.00$ & --- \\
    \bottomrule
  \end{tabular}
\end{table}

算例 A 是常规情形：目标功率高于空载损耗，正根有意义，
回代后功率为 $39.99\ \mathrm{W}$，与目标 $40\ \mathrm{W}$ 相差 $0.03\%$，
误差来自计算过程中的舍入。负根高达 $-9.99\times10^4$，
远超 C620 的控制量范围 $[-16384,16384]$，因此实际上根本不可达。
这说明对高转速工况而言，负根在\emph{数学上}存在、在\emph{执行器上}却不可达。

算例 B 和 C 暴露了 \fref{eq:rootselect} 的两个失效场景，下一节详细讨论。

\begin{figure}[htbp]
  \centering
  \begin{tikzpicture}
    \begin{axis}[
      width=0.88\textwidth, height=6.6cm,
      xlabel={转子转速 $\omega$ (rpm)},
      ylabel={允许的发送值上限},
      xmin=0, xmax=9000, ymin=0, ymax=19500,
      scaled y ticks=false,
      ytick={0,4000,8000,12000,16000},
      grid=both, grid style={gray!18},
      legend style={at={(0.97,0.97)}, anchor=north east, draw=gray!40,
                    fill=white, fill opacity=0.9, text opacity=1, font=\small},
      tick label style={font=\small},
      label style={font=\small},
      axis line style={gray!60}
    ]
      \addplot[thick, codegreen, dashed] coordinates {
        (0,18240) (500,14619) (1000,11817) (1500,9693) (2000,8086) (2500,6856)
        (3000,5898) (4000,4517) (5000,3577) (6000,2892) (7000,2369) (8000,1950) (9000,1604)};
      \addlegendentry{模型解（未受执行器限制）}
      \addplot[thick, eblue] coordinates {
        (0,16384) (500,14619) (1000,11817) (1500,9693) (2000,8086) (2500,6856)
        (3000,5898) (4000,4517) (5000,3577) (6000,2892) (7000,2369) (8000,1950) (9000,1604)};
      \addlegendentry{恒功率包络（实际可发送上限）}
      \addplot[gray, thick, dotted] coordinates {(0,10000) (9000,10000)};
      \addlegendentry{某一固定力矩阈值（示意）}
      \addplot[red!70!black, thick, dashdotted] coordinates {(0,16384) (9000,16384)};
      \addlegendentry{C620 控制量上限}
    \end{axis}
  \end{tikzpicture}
  \caption{单台电机在功率预算 $45\ \mathrm{W}$ 下的恒功率包络。
  固定力矩阈值在低速段偏松、在高速段偏紧，而恒功率包络随转速单调下降；
  两者只在某一个转速附近偶然相交。}
  \label{fig:envelope}
\end{figure}

\subsection{两类边界情形及其修复}

\paragraph{情形 I：目标功率低于空载损耗，正根变为负数。}
算例 B 中 $P^{\ast}=4\ \mathrm{W} < P_0=4.226\ \mathrm{W}$，
此时“正号根” $x_+=-114$ 是\emph{负数}。
若原始控制量为正，程序就会输出一个负的控制量，即\textbf{力矩方向被反转}。
表现在实车上就是：驾驶员在加速，底盘却给了一脚轻微刹车。

\paragraph{情形 II：判别式为负，目标功率不可达。}
算例 C 中 $P^{\ast}=2\ \mathrm{W} < P_{\min}=2.807\ \mathrm{W}$，方程无实根。

\paragraph{统一的修复方案。}
两类情形的根源是同一个：\textbf{目标功率落在了 $P_0(\omega)$ 以下，
而电机在这个转速下不可能靠正向驱动消耗这么少的功率}。
因此可以给出一个简洁而统一的处理规则：
\begin{equation}\label{eq:clip}
  P^{\ast} \;\longleftarrow\; \max\bigl(P^{\ast},\; P_0(\omega)\bigr)
  \quad\text{之后再解方程} .
\end{equation}
当 $P^{\ast}=P_0(\omega)$ 时 $C=0$，正根恰好为 $0$——即
“不发力，让电机自由转动”，这正是物理上唯一正确的选择。
\fref{eq:clip} 的好处是它同时消灭了情形 I 与情形 II，
而且代价可以忽略：被抬高的那部分功率最多就是几瓦的空载损耗，本来也无法避免。

对比之下，原实现在判别式为负时选择“跳过、维持原指令”，
这会把一个已经超功率的指令原样放行；而在情形 I 下会输出反向力矩。
两者都不是致命的（出现概率低、量值小），但在追求平顺性的底盘上是值得修掉的。

\paragraph{一个额外的数值技巧。}
当 $\Delta$ 接近零、或 $|B|$ 与 $\sqrt{\Delta}$ 量级接近时，
\fref{eq:roots} 中会出现“两个相近的大数相减”，有效数字被严重抵消。
数值分析中的标准做法~\cite{press2007numerical} 是先算
\begin{equation}\label{eq:stabq}
  q = -\frac{1}{2}\Bigl(B + \operatorname{sgn}(B)\sqrt{\Delta}\Bigr),
\end{equation}
再由 $x_1 = q/A$、$x_2 = C/q$ 得到两个根。
其中 $x_2$ 借助乘积关系得到，避开了相减抵消。
在 STM32 上用单精度浮点计算时，这个技巧能明显改善临界工况下的精度。
对本文这个方程，$B$ 大约在 $10^{-3}\sim10^{-2}$ 量级、
$\sqrt\Delta$ 与之同量级，属于\emph{恰好}容易抵消的区间，因此并非可有可无的优化。

把这一节的完整过程整理成流程图，就是图~\ref{fig:inverse}。
它是整个算法里唯一需要开方的地方，也是唯一边界情形会真正出现的地方。

\begin{figure}[htbp]
  \centering
  \begin{tikzpicture}[>=Stealth, node distance=3.2mm]
    \node[fio] (in) at (0,0) {目标功率 $P^{\ast}$、转速 $\omega$、原始指令 $\Icmd$};

    \node[fblock] (c0) at (0,-1.7)
      {算空载损耗 $P_0 = k_2\omega^2 + a$\\[1pt]
       \scriptsize 钳位 $P^{\ast}\gets\max\bigl(P^{\ast},\,P_0\bigr)$，避免出现反向力矩};

    \node[fblock] (c1) at (0,-3.7)
      {$B = \CT\omega$，\quad $C = P_0 - P^{\ast}$\\[1pt]
       $\Delta = B^2 - 4k_1C$};

    \node[draw=eblue, thick, diamond, aspect=2.4, inner sep=1pt,
          fill=eblue!6, align=center, font=\small] (d) at (0,-5.5) {$\Delta\ge0$ ?};

    \node[fblock] (yes) at (0,-7.4)
      {按原始指令 $\Icmd$ 的符号选根\\[1pt]
       $\Icmd' = \dfrac{-B \pm \sqrt{\Delta}}{2k_1}$};

    \node[fwarn] (no) at (6.2,-5.5) {目标不可达\\[1pt]\scriptsize 置 $\Icmd'=0$};

    \node[fblock] (clamp) at (0,-9.3) {限幅到 $[-16384,\;16384]$};
    \node[fio] (out) at (0,-10.8) {输出 $\Icmd'$};

    \draw[farr] (in) -- (c0);
    \draw[farr] (c0) -- (c1);
    \draw[farr] (c1) -- (d);
    \draw[farr] (d) -- node[flab, anchor=west]{是} (yes);
    \draw[farr] (d) -- node[flab, anchor=south]{否} (no);
    \draw[farr] (yes) -- (clamp);
    \draw[farr] (no.south) |- (clamp.east);
    \draw[farr] (clamp) -- (out);
  \end{tikzpicture}
  \caption{给定功率反解指令的实现流程。左侧那条“钳位”是最容易被忽略、
  却最容易被实车遇到的一步：目标功率低于空载损耗时，
  不钳位会让解出的指令变成负号，也就是力矩反向。}
  \label{fig:inverse}
\end{figure}

% ============================================================================
\section{多电机功率再分配}\label{sec:redistribute}

前面解决了单台电机的问题。真正的底盘上有四台电机，
需要在“总功率不超过预算”这一约束下，把有限的功率分给四台电机。
这一步是本文所讨论的算法与以往方案分野的地方。

\subsection{问题的提法}

设四台电机的原始控制量为 ${\Icmd}_i$、转速为 $\omega_i$（$i=1,\dots,4$），
当前控制周期可用的总功率预算为 $\Plim$。要求：
\begin{enumerate}
  \item $\sum_{i=1}^{4} P_i \le \Plim$（硬约束）；
  \item 尽可能保留驾驶员的原始意图（各轮之间的相对关系），
  否则底盘会失去方向性、走不直；
  \item 在满足前两条的前提下尽可能把预算用满，以获得最大动力。
\end{enumerate}

\subsection{方案一：功率均分}

最直接的想法是每个电机分 $\Plim/4$，各自独立反解。这个方案很快就被排除，原因有二：

\begin{itemize}
  \item \textbf{与转速耦合，破坏方向性。} 由 图~\ref{fig:powercurve} 可见，
  同样的功率预算在不同转速下对应完全不同的控制量。
  四个轮子只要转速略有差异，得到的力矩上限就不同，
  而力矩差又会产生额外的偏航力矩，进一步加剧转速差，形成正反馈。
  直线行驶时表现为“跑偏”，转向时会与驾驶意图打架。
  \item \textbf{预算无法转移。} 若某个电机本来就用不到它那一份预算
  （例如对角轮载荷很小），剩余预算不能被其他电机使用，白白浪费。
\end{itemize}

\subsection{方案二：取绝对值后统一缩放}

广工方案的做法是：把模型中的各量取绝对值，计算出“总消耗功率”，
再求一个统一缩放系数 $k$ 去缩放力矩。这个方案比均分好得多，
因为它保留了各轮之间的相对比例，方向性得到保证。但它有两个结构性缺陷。

\paragraph{缺陷一：丢失符号，把回馈制动误判为耗电。}
当一个电机的力矩与转速反向时（减速工况），它在向母线回馈能量，实际功率为负。
取绝对值之后，这一项变成一个\emph{正的}大功率，被计入总消耗。
于是控制器以为底盘比实际耗电更多，从而过度收缩力矩。用开源报告的话说：

\begin{quote}
\itshape
“因为在负功情况下根本不消耗功率，但是使用该式子计算依然会得到一个较大的‘正’功率，
这会导致对于功率利用的严重不充分。”
\end{quote}

\paragraph{缺陷二：力矩的等比缩放不等于功率的等比缩放。}
这是更本质的问题。设某电机的原始功率为 $P$，力矩缩放系数为 $k$，则缩放后的功率为
\begin{align}
  P(k\,\Icmd) &= k_2\omega^2 + k\,\CT\Icmd\omega + k^{2}k_1\Icmd^{2} + a ,\\
  k\,P(\Icmd) &= k\,k_2\omega^2 + k\,\CT\Icmd\omega + k\,k_1\Icmd^{2} + k\,a .
\end{align}
两者之差为
\begin{equation}\label{eq:scalegap}
  P(k\Icmd) - k\,P(\Icmd)
  = (1-k)\Bigl[k_2\omega^2 + a - k\,k_1\Icmd^{2}\Bigr] .
\end{equation}
只要 $k\ne1$，这个差值就一般不为零，而且它随 $\omega$、$\Icmd$ 变化。
因此\textbf{事后得到的实际总功率并不等于预算 $\Plim$，偏差大小取决于工况}：
低速高力矩时 $\Icmd^2$ 项占优，高速时 $\omega^2$ 项占优，误差偏向两侧。
结果是功率控制精度随工况漂移，只能靠外层反馈去兜底。

\subsection{方案三：功率再分配（本文方案）}

本文方案的思路是把缩放做在\emph{功率}上而不是力矩上：

\begin{enumerate}
  \item \textbf{正向计算}：对每台电机，用模型算出它当前的输入功率
  \begin{equation}\label{eq:forward}
    P_i = k_2\omega_i^{2} + \CT{\Icmd}_i\omega_i + k_1{\Icmd}_i^{2} + a ,
  \end{equation}
  并求和 $P_{\mathrm{tot}} = \sum_i P_i$。注意这里的 $P_i$ \emph{保留符号}。
  \item \textbf{判断是否需要干预}：若 $P_{\mathrm{tot}} \le \Plim$，说明没有超限，
  一台电机也不动，直接返回。这一点很重要——不干预意味着不引入任何模型误差。
  \item \textbf{计算功率缩放系数}
  \begin{equation}\label{eq:kscale}
    k = \frac{\Plim}{P_{\mathrm{tot}}} \qquad (0<k<1).
  \end{equation}
  \item \textbf{反解}：对每台电机令目标功率 $P_i^{\ast} = kP_i$，
  代入第~\ref{sec:inverse} 节的方法求出新的控制量 ${\Icmd}_i'$，
  并按原始控制量的符号选择根。
\end{enumerate}

这个方案有三个值得强调的性质。

\paragraph{性质一：精确命中预算。}
若四台电机都参与缩放，则
\begin{equation}\label{eq:exact}
  \sum_i P_i^{\ast} = \sum_i kP_i = k\sum_i P_i = \frac{\Plim}{P_{\mathrm{tot}}}\cdot P_{\mathrm{tot}} = \Plim .
\end{equation}
也就是说，\emph{目标功率之和}恰好等于预算。
在模型准确的假设下，实际总功率就等于预算——这是力矩缩放方案做不到的。
\fref{eq:scalegap} 与 \fref{eq:exact} 的对比，正是“功率再分配”这个名字的由来。

\paragraph{性质二：按消耗比例分担削减。}
每台电机的功率按同一个比例 $k$ 削减，
因此“本来消耗多的电机削减得多、本来消耗少的削减得少”，
各轮之间的功率关系保持等比。这比均分公平，也比固定力矩阈值合理。

\paragraph{性质三：负功率电机不缩放。}
对于一个正在回馈制动的电机（$P_i<0$），把它再缩小没有意义——
它本来就没有消耗功率，缩放它只会削弱制动能力。
因此实现中直接跳过它，保留原始控制量。这还顺带带来一点收益：
负功率项在 $P_{\mathrm{tot}}$ 中本来就以负号参与，
不缩放它相当于让其他电机多分摊一点，总功率会更接近而不超过预算。
下一小节给出严格的代数分析。

\subsection{负功率电机的处理：一个容易写错的细节}

设正向消耗的电机功率之和为 $S>0$，回馈电机的功率之和为 $-N$（$N>0$），
因此 $P_{\mathrm{tot}} = S - N$。我们只在 $S-N>\Plim$ 时才缩放。

\paragraph{实现 A：负功率计入分母（等价于代码的实际行为）。}
$k = \Plim/(S-N)$。缩放后正向电机总功率为 $kS$，回馈电机不变，于是实际总功率
\begin{equation}\label{eq:caseA}
  P_{\mathrm{tot}}^{A} = \frac{\Plim\,S}{S-N} - N
  = \Plim + \underbrace{\Bigl(\frac{\Plim\,N}{S-N} - N\Bigr)}_{\text{偏差}} .
\end{equation}
由于进入缩放分支意味着 $S-N>\Plim$，故 $\Plim/(S-N)<1$，偏差项为负，即
$P_{\mathrm{tot}}^{A} \le \Plim$。
所以这种做法是\emph{安全}的，只是会略微少用一点预算。

\paragraph{实现 B：负功率排除在分母之外。}
$k = \Plim/S$，缩放后正向电机总功率恰为 $\Plim$，实际总功率
\begin{equation}\label{eq:caseB}
  P_{\mathrm{tot}}^{B} = \Plim - N .
\end{equation}
两种做法的差为
\begin{equation}
  P_{\mathrm{tot}}^{A} - P_{\mathrm{tot}}^{B}
  = \Plim\Bigl(\frac{S}{S-N}-1\Bigr) - N + N
  = \frac{\Plim\,N}{S-N} > 0 ,
\end{equation}
即实现 A 用掉的功率更多、更接近预算。
这与开源报告中的说法一致——“该方法的好处是不光可以减少输出负功的电机功率带来的误差，
还可以相应的使总输出功率略微变大，尝试令其他电机吸收这部分负功”。

\paragraph{一处需要留意的实现细节。}
原代码中这一句是
\begin{lstlisting}[style=cstyle, numbers=none]
float initial_give_power[4];
...
    initial_give_power[i] = ...;
    if (initial_give_power < 0)   // negative power not included (transitory)
        continue;
    initial_total_power += initial_give_power[i];
\end{lstlisting}
\code{initial\_give\_power} 是数组名，在表达式里会退化为指针，
因此 \code{initial\_give\_power < 0} 是\textbf{指针与零的比较}，恒为假。
结果是：负功率电机实际上\emph{始终被计入}了总和，
真正生效的是上面的“实现 A”。
若要严格实现注释所描述的语义，应当写成 \code{initial\_give\_power[i] < 0}。
好在由前面的推导可知实现 A 是安全的，所以这处笔误没有造成超功率，
只是让实际行为与注释不一致，给调试带来困惑。

整个分配环节的实现流程如图~\ref{fig:redistribute} 所示。
它由一次“正向遍历”、一次判断、一次求 $k$、再一次“反向遍历”构成，
两次遍历都是对四个电机各算一遍，计算量在一次控制周期内微不足道。

\begin{figure}[htbp]
  \centering
  \begin{tikzpicture}[>=Stealth, node distance=3.2mm]
    \node[fio] (in) at (0,0) {四个电机各自的 ${\Icmd}_i$ 与 $\omega_i$};

    \node[fblock] (fwd) at (0,-1.7)
      {正向预测：$P_i = k_2\omega_i^2 + \CT{\Icmd}_i\omega_i + k_1{\Icmd}_i^2 + a$\\[1pt]
       \scriptsize 逐台计算，\textbf{保留符号}};
    \node[fblock] (tot) at (0,-3.5) {求和 $P_{\mathrm{tot}} = \sum_i P_i$};

    \node[draw=eblue, thick, diamond, aspect=2.4, inner sep=1pt,
          fill=eblue!6, align=center, font=\small] (d) at (0,-5.3)
      {$P_{\mathrm{tot}} > \Plim$ ?};

    \node[fio] (skip) at (6.3,-5.3) {不干预\\[1pt]\scriptsize 四路指令原样下发};

    \node[fblock] (k) at (0,-7.2) {求功率缩放系数 $k = \Plim \,/\, P_{\mathrm{tot}}$};

    \node[fblock] (neg) at (0,-9.0)
      {若 $P_i < 0$（正在回馈制动）：保持 ${\Icmd}_i' = {\Icmd}_i$};
    \node[fblock] (pos) at (0,-10.6)
      {否则：$P_i^{\ast} = \max\bigl(kP_i,\; k_2\omega_i^2+a\bigr)$，
       代入图~\ref{fig:inverse} 反解 ${\Icmd}_i'$};

    \node[fio] (out) at (0,-12.2) {输出四路新指令};

    \draw[farr] (in) -- (fwd);
    \draw[farr] (fwd) -- (tot);
    \draw[farr] (tot) -- (d);
    \draw[farr] (d) -- node[flab, anchor=south west]{否} (skip);
    \draw[farr] (d) -- node[flab, anchor=west]{是} (k);
    \draw[farr] (k) -- (neg);
    \draw[farr] (neg) -- (pos);
    \draw[farr] (pos) -- (out);
  \end{tikzpicture}
  \caption{功率再分配环节的实现流程。关键点有两个：
  正向计算时\emph{保留功率符号}（负值代表回馈制动），
  以及反解之前把目标功率抬到空载损耗之上（否则会出现反向力矩）。}
  \label{fig:redistribute}
\end{figure}

\subsection{三种方案的对比}

\begin{table}[htbp]
  \centering
  \small
  \caption{三种多电机功率分配方案对比}
  \label{tab:compare}
  \begin{tabular}{@{}>{\raggedright\arraybackslash}p{0.19\textwidth}>{\raggedright\arraybackslash}p{0.23\textwidth}>{\raggedright\arraybackslash}p{0.23\textwidth}>{\raggedright\arraybackslash}p{0.24\textwidth}@{}}
    \toprule
     & 方案一：功率均分 & 方案二：绝对值统一缩放 & 方案三：功率再分配 \\
    \midrule
    分配依据 & 每轮固定 $\Plim/4$ & 力矩等比（全局 $k$） & 功率等比（全局 $k$） \\
    是否保留功率符号 & 保留，但因均分而失效 & 不保留 & 保留 \\
    总功率是否精确命中 & 否（受负载影响） & 否（见 \fref{eq:scalegap}） & 是（见 \fref{eq:exact}） \\
    方向性（走直） & 差 & 好 & 好 \\
    回馈能量利用 & 无法利用 & 完全丢弃 & 可部分利用 \\
    计算量 & 低 & 低 & 中 \\
    主要误差来源 & 转速差异 & 符号丢失 $+$ 缩放非线性 & 模型精度 \\
    \bottomrule
  \end{tabular}
\end{table}

\subsection{算法伪代码}

\begin{algorithm}[htbp]
  \caption{基于功率再分配的底盘功率控制（单控制周期）}
  \label{alg:redistribute}
  \begin{algorithmic}[1]
    \Require 四台电机的原始控制量 ${\Icmd}_i$、转速 $\omega_i$；功率预算 $\Plim$
    \Ensure 四台电机的最终控制量 ${\Icmd}_i'$
    \State $P_{\mathrm{tot}} \gets 0$
    \For{$i = 1$ to $4$}
      \State $P_i \gets k_2\omega_i^{2} + \CT{\Icmd}_i\omega_i + k_1{\Icmd}_i^{2} + a$
      \State $P_{\mathrm{tot}} \gets P_{\mathrm{tot}} + P_i$
    \EndFor
    \If{$P_{\mathrm{tot}} \le \Plim$}
      \State \Return ${\Icmd}_i' \gets {\Icmd}_i$ \Comment{未超限，不干预}
    \EndIf
    \State $k \gets \Plim / P_{\mathrm{tot}}$
    \For{$i = 1$ to $4$}
      \If{$P_i < 0$}
        \State ${\Icmd}_i' \gets {\Icmd}_i$ \Comment{回馈制动，不缩放}
      \Else
        \State $P_i^{\ast} \gets \max\bigl(kP_i,\; k_2\omega_i^{2}+a\bigr)$
        \Comment{式~\eqref{eq:clip}}
        \State $B \gets \CT\omega_i$;\quad $C \gets k_2\omega_i^{2} + a - P_i^{\ast}$
        \State $\Delta \gets B^{2} - 4k_1C$
        \If{$\Delta < 0$}
          \State ${\Icmd}_i' \gets 0$
        \ElsIf{${\Icmd}_i > 0$}
          \State ${\Icmd}_i' \gets \min\bigl((-B+\sqrt{\Delta})/(2k_1),\; 16384\bigr)$
        \Else
          \State ${\Icmd}_i' \gets \max\bigl((-B-\sqrt{\Delta})/(2k_1),\; -16384\bigr)$
        \EndIf
      \EndIf
    \EndFor
  \end{algorithmic}
\end{algorithm}

\subsection{两个可选的改进方向}

\paragraph{一轮修正。}
由于 \fref{eq:clip} 的钳位以及负功率电机不缩放，
实际总功率通常会略低于预算。若希望把预算用得更满，
可以在第 5 步之后重新正向计算一次实际总功率 $P_{\mathrm{tot}}'$，
若 $P_{\mathrm{tot}}' < \Plim$ 且余量可观（例如超过 $5\%$），
再用 $k' = \Plim/P_{\mathrm{tot}}'$ 做一次补偿。
由于模型对控制量是单调的，这种迭代会快速收敛，通常一到两轮就足够。

\paragraph{按余量而非按比例分配。}
比例分配在“某一台电机已经接近其物理极限”时并不理想。
更精细的做法是把预算看作一个资源分配问题：
按各电机在当前位置上“单位功率能换来多少力矩”的边际效率来分配，
即求解一个小规模的凸优化问题。代价是计算量上升，
在 STM32F4 上需要谨慎评估。这项工作留待第~\ref{sec:future} 节讨论。

% ============================================================================
\section{缓冲能量闭环}\label{sec:bufferloop}

前面的推导给出了一个“给定功率预算 $\Plim$，把力矩压到预算之内”的开环执行机构。
但预算本身从哪里来？如果直接把裁判系统的功率上限 $\Pmax$ 当作预算，
那么模型误差、母线电压波动、绕组升温都会让实际功率偏离预算，
底盘要么长期超功率、要么长期浪费。因此还需要一层外环。
本节说明这一层外环的物理背景、控制结构与参数整定。

\subsection{缓冲能量的官方结算规则}

RoboMaster 规则手册对缓冲能量的描述可以概括为一个离散积分器。
裁判系统以 $10\ \mathrm{Hz}$（每 $100\ \mathrm{ms}$）结算一次：
\begin{equation}\label{eq:officialbuffer}
  Z \;\leftarrow\; Z - (P_r - P_l)\times 0.1,
  \qquad Z\in[0,\,60]\ \mathrm{J},
\end{equation}
其中 $P_r$ 为底盘的瞬时功率，$P_l$ 为该等级对应的上限功率。
把它写成连续时间形式\cite{rm2024rules} 就是
\begin{equation}
  \dot Z = (P_l - P_r)\times 10 \quad [\mathrm{J/s}],
\end{equation}
即\emph{低于上限时回升的速率与超过上限时下降的速率完全相同}。这一点值得特别指出：
规则手册并没有单独规定一个“恢复速率”，网上流传的所谓 $24\ \mathrm{J/s}$、
$36\ \mathrm{J/s}$ 等具体数值并无可查证的官方依据，不应作为设计输入。

此外还有两条容易被忽略的规则细节：
\begin{itemize}
  \item 触发“飞坡增益”后，缓冲能量上限临时提高到 $250\ \mathrm{J}$；
  但一旦被消耗到 $60\ \mathrm{J}$ 以下，此后最多只能恢复到 $60\ \mathrm{J}$。
  \item 缓冲能量耗尽后进入扣血阶段，扣血量按超限比例分档：
  \begin{equation}
    K = \frac{P_r - P_l}{P_l}\times100\%,
    \qquad
    N = \begin{cases}
      10\%, & K \le 10\%,\\
      20\%, & 10\% < K \le 20\%,\\
      40\%, & K > 20\%.
    \end{cases}
  \end{equation}
  每检测周期扣血 $=$ 上限血量 $\times N\% \times 0.1$。
  也就是说，超限一倍（$K=100\%$）时每秒扣掉 $40\%$ 的血量——
  这解释了为什么“宁可保守也不要超功率”。
\end{itemize}

\subsection{为什么需要闭环}

设本周期实际执行时，电机功率与指令之间存在一个偏差 $\delta$（模型误差、电压波动、
温度漂移、机械负载变化都会贡献到 $\delta$），即
\begin{equation}
  \Pin = \Plim + \delta .
\end{equation}
于是缓冲能量的动态变为
\begin{equation}\label{eq:plant}
  \dot\Ebuf = \Pmax - \Pin = \Pmax - (\Pmax - u) - \delta = u - \delta,
\end{equation}
其中 $u$ 是外环控制器的输出。\fref{eq:plant} 说明\textbf{被控对象是一个纯积分器}：
控制量 $u$（单位 W）积分后得到 $\Ebuf$（单位 J），直流增益为 $1$，
而模型误差 $\delta$ 相当于作用在积分器\emph{输入端}的常值扰动。

这个结构决定了外环的设计：

\begin{itemize}
  \item 若只用比例控制（$u = K_p e$，$e=\Ebuf^{\ast}-\Ebuf$），
  闭环为 $\dot e = -K_p e + \delta$，稳态误差为
  \begin{equation}
    e_{\mathrm{ss}} = \frac{\delta}{K_p} .
  \end{equation}
  也就是说，\emph{模型误差会直接变成一个缓存的静差}。
  由于对象本身含积分器，闭环对参考值阶跃是无静差的；
  但对“作用在对象输入端的扰动”，纯比例控制只能把误差压缩 $K_p$ 倍而不能消除——
  这是系统类型（system type）理论的标准结论。
  \item 引入积分项后，开环从 I 型升为 II 型，常值扰动被完全抑制，静差归零。
  代价是多了一个自由度需要整定，且积分饱和会带来超调。
\end{itemize}

结论是：\textbf{外环应当用 PI（或 PID）控制，且积分项不可省}。

\subsection{控制器结构与闭环分析}

取
\begin{equation}\label{eq:outerloop}
  e = \Ebuf^{\ast} - \Ebuf,
  \qquad
  u = K_p e + K_i\!\int e\,\dd t,
  \qquad
  \boxed{\;\Plim = \Pmax - u\;}
\end{equation}
把 \fref{eq:outerloop} 代入 \fref{eq:plant}：
\begin{equation}
  \dot e = -\dot\Ebuf = -u + \delta = -K_p e - K_i\!\int e\,\dd t + \delta .
\end{equation}
两边再对时间求导，$\delta$ 的导数为零（常值扰动），得到齐次方程
\begin{equation}\label{eq:closedloop}
  \boxed{\;\ddot e + K_p\,\dot e + K_i\,e = 0\;}
\end{equation}
其闭环特征方程为
\begin{equation}\label{eq:chareq}
  s^{2} + K_p s + K_i = 0 ,
\end{equation}
与标准二阶系统 $s^2 + 2\zeta\omega_n s + \omega_n^2 = 0$ 对照（见文献~\cite{franklin2002feedback}），立刻得到整定公式：
\begin{equation}\label{eq:tuning}
  \boxed{\;\omega_n = \sqrt{K_i}, \qquad \zeta = \frac{K_p}{2\sqrt{K_i}}\;}
  \quad\Longleftrightarrow\quad
  K_p = 2\zeta\omega_n,\quad K_i = \omega_n^{2}.
\end{equation}

\fref{eq:chareq} 还澄清了一个常见误解：这个环是\emph{二阶}的，
它的稳定性与响应速度完全由 $K_p,K_i$ 决定，与电机模型、底盘质量无关。
这一点让整定变得非常可控。

\begin{note}[关于积分项与采样周期]
很多 RM 工程代码里的 PID 实现是
\code{integral += err; out = Kp*err + Ki*integral;}，
即积分项\emph{不显式乘以 $\dd t$}。在这种约定下，
代码里的 $K_i^{\mathrm{code}}$ 与 \fref{eq:tuning} 中的连续域 $K_i$ 相差一个采样周期：
\begin{equation}
  K_i^{\mathrm{code}} = K_i\cdot T_s .
\end{equation}
因此\textbf{一旦改变了调用周期 $T_s$，必须同步重算 $K_i$}，
否则同样的参数会得到完全不同的动态。
\end{note}

\begin{table}[htbp]
  \centering
  \small
  \caption{缓冲能量闭环整定示例（$T_s = 1\ \mathrm{ms}$，积分项不乘 $\dd t$）}
  \label{tab:tuning}
  \begin{tabular}{@{}cccccc@{}}
    \toprule
    设计目标 & $\zeta$ & $\omega_n$ (rad/s) & $K_p$ & $K_i$（连续） & $K_i^{\mathrm{code}}$ \\
    \midrule
    缓慢、无超调 & $1.0$ & $1.0$  & $2.0$  & $1.0$  & $0.001$ \\
    折中（推荐） & $0.8$ & $3.0$  & $4.8$  & $9.0$  & $0.009$ \\
    较快 & $0.7$ & $6.0$  & $8.4$  & $36.0$ & $0.036$ \\
    \bottomrule
  \end{tabular}
\end{table}

带宽的上限受两个因素约束：裁判系统的结算频率（$10\ \mathrm{Hz}$）
与执行机构的更新频率。把闭环带宽取到结算频率的 $1/10$ 以下（约 $1\ \mathrm{Hz}$ 以内）
是比较稳妥的做法，\fref{eq:tuning} 中 $\omega_n=3\ \mathrm{rad/s}$（约 $0.5\ \mathrm{Hz}$）
就落在这个区间里。

\subsection{极性：一个必须核对的问题}

\fref{eq:outerloop} 中的减号不是可有可无的记号。它的物理含义是：

\begin{quote}
当缓冲能量\emph{偏低}（$e>0$）时，控制器必须\emph{收紧}功率预算（$\Plim<\Pmax$），
让底盘少用一点功率，缓冲能量才能回充；反之当缓冲能量偏高时才放宽预算。
\end{quote}

如果极性写反，闭环变成正反馈：缓冲能量低 $\to$ 放宽预算 $\to$ 更超功率 $\to$
缓冲能量更低，很快触发扣血。
\begin{equation}
  \text{判据：}\quad
  \frac{\partial \Plim}{\partial \Ebuf} > 0 .
\end{equation}
即“缓冲能量越多，允许的功率预算越大”。任何实现都应先按这条判据自检一遍。

\begin{remark}
开源报告中的控制框图把 PID 输出到求和点的符号画成了正号，
而对应的代码写的是 \code{input\_power = max\_power\_limit - buffer\_pid.out}，
两者在符号上是不一致的。按照上面的判据，\textbf{代码是正确的一方}。
在阅读任何一份功率控制框图时，都建议先用这条判据验一遍极性。
\end{remark}

\subsection{工程上的若干改进}

\paragraph{(a) 用 $\sqrt{\cdot}$ 误差映射代替线性误差。}
一个被多个队伍采用的技巧~\cite{hkustpowermodule} 是把误差定义在“能量平方根”上：
\begin{equation}\label{eq:sqrterror}
  e = \sqrt{\Ebuf^{\ast}} - \sqrt{\Ebuf},
  \qquad
  \Plim = \Pmax - K_p e - K_d\frac{e - e_{\mathrm{last}}}{\Delta t}.
\end{equation}
由于 $\dfrac{\dd\sqrt{E}}{\dd E} = \dfrac{1}{2\sqrt{E}}$，
同一个能量偏差在能量较高时对应更小的 $e$，因此这是一个\emph{增益随能量上升而减小}的
非线性映射：在接近满电时环路更柔和，避免功率上限在 $E$ 靠近上限时剧烈振荡；
在能量偏低时增益更大，能更快地收紧预算。
实现上只是把两次开方换个位置，代价极小，收益明显。

\paragraph{(b) 微分项要慎用。}
微分项有助于抑制超调，但由于缓冲能量的更新是 $10\ \mathrm{Hz}$ 的阶梯信号，
直接差分会放大量化噪声，实际效果往往是引入振荡。
工程经验是“宁可不用 D，也不要调不好的 D”；若一定要用，
必须先对 $\Ebuf$ 做低通滤波再差分。

\paragraph{(c) 积分饱和与抗饱和。}
$\Plim$ 不能为负，$\Ebuf$ 也在 $[0,E_{\max}]$ 内被限幅，
因此积分器会在限幅期间持续累积（windup），造成解除限幅后的大幅超调。两条对策：
\begin{itemize}
  \item 对积分项显式限幅：$|u_I| \le u_I^{\max}$；
  \item 对输出做条件积分（conditional integration）：当 $\Plim$ 已经触到 $0$ 或触到上限时，
  停止积分。
\end{itemize}

\paragraph{(d) 通信异常时的兜底。}
若在预期时间内没有收到裁判系统的缓冲能量数据，或者解析出的数值明显不合理，
应当\emph{把预算严格设为 $\Pmax$}（保守策略），而不是沿用旧值或直接放行。
这是一条来自官方圆桌~\cite{rmroundtable006} 与多个开源实现的一致建议。

\subsection{与超级电容的协同}

在超电系统中，外环输出的 $\Plim$ 不再直接等于底盘电机可以消耗的功率上限，
而是“允许底盘从电池侧取用的功率”。超电控制器负责补足差额，
因此主控还需要把这一信息通过 CAN 告知超电控制器。
开源实现中的做法是：

\begin{lstlisting}[style=cstyle, numbers=none]
input_power = max_power_limit - chassis_power_control->buffer_pid.out;
CAN_CMD_CAP(input_power);            // 告知超电：本周期计划从电池取多少功率
...
if (cap_measure.cap_percent > 5) {
    if (cap_state == 0) chassis_max_power = input_power + 5;    // 常规档
    else                chassis_max_power = input_power + 200;  // 爆发档
} else {
    chassis_max_power = input_power;                            // 电容电量低，不借力
}
\end{lstlisting}

这段代码里有三个值得说明的设计点。

\begin{enumerate}
  \item \textbf{常规档额外放宽 $5\ \mathrm{W}$。}
  让底盘长期“略微超出电池侧的目标功率”，从而持续小额消耗电容能量。
  这样电容不会长期处于满电状态，一旦出现回馈制动就有空间吸收能量。
  原代码注释把这一做法的目的表述为“提高能量利用率”。
  \item \textbf{爆发档额外放宽 $200\ \mathrm{W}$。}
  由遥控器按键切换到主动释放电容能量的模式。它与港科开源实现中
  “补电模式”的思路一致——把功率上限设成远高于裁判限制，
  由超电长时间补偿差额，从而在一段时间内“以低等级的功率限制打出高等级的动力”。
  需要注意的是这种模式的可持续时间完全由电容储量决定，
  且电容电压越低，可补偿的功率越小，因此必须配合电量保护。
  \item \textbf{电量保护。}
  当电容剩余电量低于 $5\%$ 时，取消一切额外放宽，回到严格的 $\Pmax$。
  这是必要的：电容电压降低后，同样的放电电流能提供的功率下降，
  继续下达超功率指令既不安全也会直接超功率。
\end{enumerate}

\begin{note}
超电接入之后，“缓冲能量由谁负责”这件事本身发生了变化。
一个被广泛接受的分工~\cite{zjuwiki} 是：\emph{缓冲能量的观测与能量调度下沉到超电控制器}，
主控只负责把功率预算分发到四个电机。
因此当出现“电容还有电却仍然超功率”的现象时，
应当优先检查超电侧的功率环与硬件，而不是继续拧主控的 PID。
第~\ref{sec:supercap} 节会对这一分层关系展开讨论。
\end{note}

\subsection{整体控制框图}

\begin{figure}[htbp]
  \centering
  \begin{tikzpicture}[
      font=\small,
      >={Stealth[length=1.8mm]},
      block/.style={draw=eblue, thick, fill=eblue!4, rounded corners=1.5pt,
                    minimum height=8mm, minimum width=14mm, align=center, inner sep=2pt},
      bigblock/.style={block, minimum width=26mm},
      sum/.style={draw=eblue, thick, circle, inner sep=0pt, minimum size=5mm, fill=white},
      line/.style={draw, thick, -{Stealth[length=1.8mm]}},
      pline/.style={draw, thick}
    ]
    \node[sum] (s1) {};
    \node[block, right=8mm of s1] (pid) {PID$(t)$};
    \node[sum, right=9mm of pid] (s2) {};
    \node[bigblock, right=9mm of s2] (gc) {功率再分配\\$G_c(t)$};
    \node[block, right=11mm of gc] (gm) {电机\\$G_{\mathrm{motor}}$};
    \node[bigblock, below=9mm of pid] (int) {缓冲能量结算\\$\int(P_{\max}-\Pin)\dd t$};

    \draw[line] (-1.9,0) -- node[above, pos=0.62]{$30$ J} node[below, pos=0.62, font=\scriptsize]{$+$} (s1);
    \draw[line] (s1) -- (pid);
    \draw[line] (pid) -- node[above, font=\scriptsize]{$-$} (s2);
    \draw[line] (s2) -- node[above]{$\Plim$} (gc);
    \draw[line] (gc) -- node[above]{$\Icmd$} (gm);
    \draw[line] (gm.east) -- ++(1.0,0) coordinate (pout);
    \node[right=1mm of pout] {$\Pin$};

    \node[above=6mm of s2] (pmax) {$\Pmax$};
    \draw[line] (pmax) -- node[right, font=\scriptsize]{$+$} (s2);

    \draw[line] (pout) |- (int.east);
    \draw[line] (int.west) -- ++(-1.0,0) coordinate (fb);
    \draw[line] (fb) |- (s1.south);
    \node[font=\scriptsize, anchor=south east] at ([xshift=-1mm, yshift=1mm]s1.south)
      {$\Ebuf$（$-$）};

    \node[draw=gray, dashed, rounded corners=3pt, inner sep=7pt,
          fit=(s1)(pid)(s2)(gc)(int), label={[gray, font=\scriptsize]above:软件部分}] {};
  \end{tikzpicture}
  \caption{底盘功率控制的完整框图。虚线内为软件部分，包含一个缓冲能量外环
  与一个功率再分配环节；$G_{\mathrm{motor}}$ 为电机真实的物理模型。
  注意 PID 输出在求和点处为\emph{负}号，即 $\Plim = \Pmax - u$，
  这与上一小节讨论的极性判据一致。}
  \label{fig:blockdiagram}
\end{figure}

% ============================================================================
\section{代码实现逐行解读}\label{sec:code}

本节对开源实现 \code{chassis\_power\_control.c} 做一次逐行解读，
指出其中若干处可能导致行为偏差的实现细节。
代码来自文献~\cite{lin2023report,lin2023code}（MIT 许可），为便于阅读，
保留了原注释并统一了缩进。

\subsection{参数与数据结构}

\begin{lstlisting}[style=cstyle, numbers=none]
uint16_t max_power_limit = 40;
fp32  toque_coefficient = 1.99688994e-6f;  // (20/16384)*(0.3)*(187/3591)/9.55
fp32  a        = 1.23e-07;                 // 文档中的 k1（I_cmd^2 的系数）
fp32  k2       = 1.453e-07;                // 文档中的 k2（omega^2 的系数）
fp32  constant = 4.081f;                   // 文档中的常数项 a
\end{lstlisting}

这里第一处需要注意的就是\emph{命名}。如第~\ref{sec:model} 节所述，
代码中被注释为 \code{k1} 的 \code{a} 实际是 $\Icmd^2$ 的系数，
而 \code{k2} 是 $\omega^2$ 的系数；同时 \code{constant} 才是公式里的常数项 $a$。
把这段代码与任何文档对照时，都应先做一次记号映射（见表~\ref{tab:params}），
否则很容易在调参时改错对象。

\code{max\_power\_limit = 40} 是裁判系统给出的功率上限，单位 W，
由 \code{get\_chassis\_max\_power()} 覆盖。

\subsection{主循环}

\begin{lstlisting}[style=cstyle, label={lst:main}]
void chassis_power_control(chassis_move_t *chassis_power_control)
{
    fp32 chassis_max_power = 0;
    float input_power = 0;        // input power from battery (referee system)
    float initial_give_power[4];
    float initial_total_power = 0;
    fp32  scaled_give_power[4];

    get_chassis_power_and_buffer(&chassis_power, &chassis_power_buffer);
    PID_calc(&chassis_power_control->buffer_pid, chassis_power_buffer, 30);
    get_chassis_max_power(&max_power_limit);
    input_power = max_power_limit - chassis_power_control->buffer_pid.out;
    CAN_CMD_CAP(input_power);     // set the input power of capacitor controller
    ...
    for (uint8_t i = 0; i < 4; i++) {        // 正向计算：模型 -> 功率
        initial_give_power[i] =
              chassis_power_control->motor_speed_pid[i].out * toque_coefficient
                  * chassis_power_control->motor_chassis[i].chassis_motor_measure->speed_rpm
            + k2 * speed_rpm * speed_rpm
            + a  * chassis_power_control->motor_speed_pid[i].out
                 * chassis_power_control->motor_speed_pid[i].out
            + constant;
        if (initial_give_power < 0)          // <-- 见下文的讨论
            continue;
        initial_total_power += initial_give_power[i];
    }

    if (initial_total_power > chassis_max_power) {   // 超限才干预
        fp32 power_scale = chassis_max_power / initial_total_power;
        for (uint8_t i = 0; i < 4; i++) {
            scaled_give_power[i] = initial_give_power[i] * power_scale;
            if (scaled_give_power[i] < 0)
                continue;
            fp32 b = toque_coefficient * speed_rpm;
            fp32 c = k2 * speed_rpm * speed_rpm - scaled_give_power[i] + constant;
            fp32 inside = b * b - 4 * a * c;
            if (inside < 0) {
                continue;                    // <-- 维持原指令
            } else if (chassis_power_control->motor_speed_pid[i].out > 0) {
                fp32 temp = (-b + sqrt(inside)) / (2 * a);
                if (temp > 16000) temp = 16000;
                chassis_power_control->motor_speed_pid[i].out = temp;
            } else {
                fp32 temp = (-b - sqrt(inside)) / (2 * a);
                if (temp < -16000) temp = -16000;
                chassis_power_control->motor_speed_pid[i].out = temp;
            }
        }
    }
}
\end{lstlisting}

整体流程与第~\ref{sec:redistribute} 节描述的完全一致，可以逐段对应：

\begin{enumerate}
  \item \code{get\_chassis\_power\_and\_buffer} 读取裁判系统的功率与缓冲能量；
  \item \code{PID\_calc(..., 30)} 以 $30\ \mathrm{J}$ 为设定值计算缓冲能量外环；
  \item \code{input\_power = max\_power\_limit - out} 得到本周期功率预算，
  并同步给超电控制器；
  \item 第一个 \code{for} 循环做\emph{正向计算}，把模型算出的功率累加；
  \item 若未超限则\emph{完全不干预}（这是很好的设计）；
  \item 否则求功率缩放系数 \code{power\_scale}；
  \item 第二个 \code{for} 循环做\emph{反解}，用一元二次方程求根公式算出新的控制量，
  并按原始控制量 \code{motor\_speed\_pid[i].out} 的符号选择 $+$ 或 $-$ 根。
\end{enumerate}

注意 \code{2*a} 中的 \code{a} 正是 $\Icmd^2$ 的系数 $k_1$，
\code{b = toque\_coefficient * speed\_rpm} 正是 $\CT\omega$，
\code{c = k2*speed\string^2 - scaled + constant} 正是 $C$。
与 \fref{eq:invquad} 完全对应。

\subsection{值得修正的实现细节}

\begin{table}[htbp]
  \centering
  \small
  \caption{代码中值得修正的实现细节}
  \label{tab:codeissues}
  \begin{tabular}{@{}>{\raggedright\arraybackslash}p{0.05\textwidth}>{\raggedright\arraybackslash}p{0.38\textwidth}>{\raggedright\arraybackslash}p{0.47\textwidth}@{}}
    \toprule
    \# & 现象 & 影响与建议 \\
    \midrule
    1 & \code{initial\_give\_power < 0} 是数组名（指针）与 $0$ 比较，恒为假 &
    负功率电机始终被计入总和。行为上属于第~\ref{sec:redistribute} 节的“实现 A”，
    是安全的，但与注释不符。应改为 \code{initial\_give\_power[i] < 0} \\[3pt]
    2 & 变量命名与文档公式不一致（\code{a}、\code{k2}、\code{constant}） &
    极易在调参时改错对象。建议重命名为 \code{k\_omega2}、\code{k\_i2}、\code{c0} \\[3pt]
    3 & 判别式为负时 \code{continue}，保留原指令 &
    原指令可能已超功率。建议按 \fref{eq:clip} 把目标功率抬到 $P_0(\omega)$ 后重解，
    此时必有解且方向正确 \\[3pt]
    4 & 正根分支只做上限钳位（\code{temp > 16000}），无下限保护 &
    当目标功率低于空载损耗时正根为负，会导致\emph{力矩反向}。
    应在取根后校验符号：若与原始指令异号，则置 $0$ \\[3pt]
    5 & \code{toque\_coefficient} 拼写错误（应为 \code{torque}） &
    可读性问题，不影响功能 \\[3pt]
    6 & \code{float} 与 \code{fp32} 混用 &
    在 STM32F4 上二者等价，但混用会掩盖移植到 double 平台时的潜在差异 \\[3pt]
    7 & 每周期 4 次 \code{sqrt} 与 4 次除法 &
    1 kHz 下开销可接受；若紧张可用 \code{arm\_sqrt\_f32}，
    或把 $1/(2k_1)$ 预先算成常数 \\[3pt]
    8 & \code{chassis\_max\_power} 这一命名 &
    其实际语义是“本周期的功率预算”，而非“底盘的最大功率”。
    建议改名以免与外层的等级功率上限混淆 \\[3pt]
    9 & 母线电压反馈未被使用（虽已 \code{include} \code{voltage\_task.h}） &
    模型未做电压归一化，电池放电后段会有系统偏差。见第~\ref{sec:future} 节 \\[3pt]
    10 & 未对 $\delta$ 做任何在线修正 &
    依赖外环兜底，符合设计意图；若希望更精确，可引入第~\ref{sec:future} 节的 RLS 在线辨识 \\
    \bottomrule
  \end{tabular}
\end{table}

其中第 3、4 两条是功能性问题，其余主要是可维护性问题。
需要强调的是：这些问题的量值都不大，原实现在比赛中是可用且有效的；
这里把它们列出来，是因为它们正好对应本文前面推导出的两个边界情形，
理解它们比记住结论更有价值。

% ============================================================================
\section{工程实践：标定、部署与故障排查}

理论推导完成之后，真正决定成败的是标定质量与部署细节。
本节给出一份可执行的流程清单。

\subsection{硬件准备：电流采样量程的选择}

模型标定依赖一路可信的功率测量。常用的方案~\cite{ina226} 是在电调输入端串入分流电阻，
用 INA226 之类的双向电流/功率监测芯片测量母线电流与电压。
选型时最需要先算清楚的是\emph{量程}，因为它直接决定了数据里会不会出现削顶。

INA226 的分流电压输入满量程为 $\pm81.92\ \mathrm{mV}$（对应十进制 $7\mathrm{FFF}$），
最低有效位约 $2.5\ \mu\mathrm{V}$；总线电压输入范围为 $0\sim36\ \mathrm{V}$，
不可超过。由此可算得电流量程
\begin{equation}\label{eq:inacur}
  I_{\max} = \frac{81.92\ \mathrm{mV}}{R_{\mathrm{shunt}}} :
  \qquad
  \begin{aligned}
    R &= 2\ \mathrm{m}\Omega &&\Rightarrow I_{\max} = \pm40.96\ \mathrm{A},\\
    R &= 4\ \mathrm{m}\Omega &&\Rightarrow I_{\max} = \pm20.48\ \mathrm{A},\\
    R &= 5\ \mathrm{m}\Omega &&\Rightarrow I_{\max} = \pm16.38\ \mathrm{A}.
  \end{aligned}
\end{equation}
M3508 的额定力矩电流为 $20\ \mathrm{A}$，因此 $4\ \mathrm{m}\Omega$ 恰好覆盖，
$2\ \mathrm{m}\Omega$ 则留出一倍余量但牺牲一半分辨率。若为了压低噪声选了过大的电阻，
就必然在峰值力矩下削顶——开源报告图 1 的注记
“输入功率出现失真是因为电流超过 INA226 最大采样电流”正是这个现象。
\textbf{被削顶的样本必须剔除，否则会把最小二乘解整体拉偏}（见第~\ref{sec:lsq} 节）。

\subsection{标定流程}

\begin{enumerate}
  \item \textbf{固定硬件。} 标定与装车应使用同批、同保养状态的电机；
  齿轮箱阻尼差异会直接体现在 $k_2$ 上。若条件允许，逐台标定最优。
  \item \textbf{设计激励信号。} 用 C 型开发板向 C620 发送\emph{不同频率与幅值}的正弦力矩指令，
  同时对电机施加\emph{随机大小}的负载。
  这一步的目的是让 $\omega^2$ 与 $\Icmd^2$ 两列去相关——
  若只在恒定转速下实验，$k_2$ 与 $a$ 将不可辨识（第~\ref{sec:lsq} 节）。
  \item \textbf{采集数据。} 用 INA226 采样实时电流与母线电压，
  用 STM32CubeMonitor 把真实电流、力矩控制值、转子转速导出为 CSV。
  \item \textbf{拟合与交叉验证。} 用 MATLAB 的 \code{fit}/\code{fittype}
  或自己实现的正规方程求解（附录~\ref{app:lsq}）得到 $k_2,k_1,a$，
  再用\emph{不同批次}的数据做验证，观察参数是否稳定、验证集 RMSE 是否接近训练集。
  \item \textbf{物理自检。} 用 \fref{eq:k1phys} 反解等效相电阻，
  用 $k_2$ 估算空载损耗曲线，看二者是否落在合理范围。不合理就回到第 2 步。
  \item \textbf{装车验证。} 见第~\ref{sec:verify} 节。
\end{enumerate}

\begin{note}[三条来自开源实现的经验]
\begin{enumerate}
  \item \textbf{拟合时电机不要空载。}空载时力矩电流控制值并不等于实际力矩电流，
  因为此时电流主要被用来克服空载损耗，模型的前提被破坏。
  \item \textbf{走不直时先把功率均分试试。}这是一个极好的隔离手段：
  若均分之后仍然走不直，问题多半在模型或机械；若均分之后正常了，
  问题就在分配算法。
  \item \textbf{改完轮电机的 PID 参数后，务必同步复核功率模型的相关参数。}
  功率模型只对“最终发送值”负责；任何在功率分配之后、CAN 发送之前
  额外加入的限幅、滤波或非线性环节，都会让模型预测与实际脱节。
  最稳妥的架构是把功率控制放在整条指令链路的末端。
\end{enumerate}
\end{note}

\subsection{验证方法}\label{sec:verify}

\begin{description}
  \item[恒功率启动测试。] 让底盘从静止满油起步，同时记录四个电机的功率之和。
  若算法正确，总功率曲线应当迅速爬升到预算附近并\emph{保持平坦}，
  而不是先冲高再回落。这正是“恒功率运行”这一设计目标的直接体现。
  \item[缓冲能量稳定性测试。] 在多种工况（直线加速、原地转向、爬坡、急停）下
  观察缓冲能量是否稳定在设定值附近小幅波动。
  持续单边漂移说明外环极性或积分项有问题；周期性振荡说明带宽过高。
  \item[阶跃响应测试。] 人为把功率预算阶跃降低（例如在代码里临时改成 $20\ \mathrm{W}$），
  观察电机指令是否平滑跟随、有无过冲与振荡。
  用这条测试可以直接标定外环的实际闭环带宽。
  \item[对表测试。] 用外部功率计（或超电控制器回传的功率）与模型预测值对比，
  统计不同工况下的相对误差。误差随转速单调变化通常意味着缺项（例如电压归一化）。
\end{description}

\subsection{常见故障与排查}

\begin{table}[htbp]
  \centering
  \small
  \caption{常见故障与排查路径}
  \label{tab:trouble}
  \begin{tabular}{@{}>{\raggedright\arraybackslash}p{0.20\textwidth}>{\raggedright\arraybackslash}p{0.34\textwidth}>{\raggedright\arraybackslash}p{0.36\textwidth}@{}}
    \toprule
    现象 & 可能原因 & 排查手段 \\
    \midrule
    起步无力、加速缓慢 &
    仍在用固定力矩阈值；或 $k_1,k_2$ 明显偏大 &
    打印单电机“允许控制量上限—转速”曲线，与图~\ref{fig:envelope} 对照 \\[4pt]
    走不直、行进中偏航 &
    分配算法破坏了各轮的相对关系；或四台电机参数不一致 &
    临时改为功率均分，看是否恢复正常；再逐台标定参数 \\[4pt]
    持续超功率（缓冲能量一直被扣） &
    模型偏乐观（$k$ 偏小）；或下发了超出执行器能力的指令 &
    检查 $k_1$ 反解出的电阻是否过小；检查是否有 $|\Icmd|>16384$ 的非法值 \\[4pt]
    功率抖动、“一顿一顿” &
    外环带宽过高；或对 $10\ \mathrm{Hz}$ 的阶梯信号做了微分 &
    降低 $K_p$；去掉或前置滤波 $D$ 项 \\[4pt]
    电容还有电却超功率 &
    超电侧功率环或硬件能力不足 &
    按第~\ref{sec:bufferloop} 节的讨论，优先排查超电侧，而非继续调主控 PID \\[4pt]
    制动时超功率 &
    回馈制动功率被模型低估（符号处理不当） &
    检查建模与反解过程中是否出现过取绝对值 \\[4pt]
    通信中断后行为异常 &
    未做兜底 &
    超时未收到裁判数据时，把预算严格设为 $\Pmax$ \\
    \bottomrule
  \end{tabular}
\end{table}

% ============================================================================
\section{局限性与未来工作}\label{sec:future}

\subsection{已知局限}

\begin{enumerate}
  \item \textbf{$k_2$ 是电机的个体参数，不可通用。}
  它主要来自齿轮箱阻尼，不同电机不同。当四台电机的 $k_2$ 相差较大时，
  “按功率等比削减”会在四轮之间引入稳态误差，表现为轻微跑偏。
  \item \textbf{模型不含母线电压。} 电池从满电放到欠压，同样控制量对应的功率会下降。
  由于外环会把预算拉回平衡点，这个偏差不会导致超功率，
  但会造成“电压低时动力衰减比预期更明显”。
  \item \textbf{模型不含温度。} 铜电阻温度系数约 $+0.39\%/\mathrm{K}$，
  绕组从 $25\,^{\circ}\mathrm{C}$ 升到 $100\,^{\circ}\mathrm{C}$ 时 $k_1$ 增大三成，
  同样会被外环掩盖成“静差”，但瞬时精度会下降。
  \item \textbf{与平衡步兵的冲突。} 功率限制过严时，平衡步兵在高速下
  无法及时输出足够的反向力矩改变倾角来减速，会出现“滑铲”。
  这不是本算法的缺陷，而是“严格功率约束”与“姿态控制的瞬时功率需求”之间
  的固有矛盾。可行的方向是在检测到减速意图时临时提高功率上限，
  由超电介入提供这部分功率。
  \item \textbf{分配策略仍是“等比”而非“最优”。}
  等比削减保证了一致性，但并非在所有工况下都能最大化有效牵引力。
\end{enumerate}

\subsection{改进方向}

\paragraph{(a) 母线电压归一化。}
把 \fref{eq:model} 中的损耗项与机械功率项按 $(U_{\mathrm{nom}}/U_{\mathrm{bus}})^{2}$
与 $U_{\mathrm{nom}}/U_{\mathrm{bus}}$ 做尺度修正，等效于把模型搬到额定电压下计算。
实现成本很低（一次除法），却能在电池放电后段明显改善精度。

\paragraph{(b) 在线辨识。}
若硬件允许读取真实功率（例如超电控制器回传的母线功率），
可以用递推最小二乘（RLS）\cite{ljung1999system,hkustpowermodule} 在线更新参数：
\begin{align}
  K_k &= \frac{P_{k-1}\varphi_k}{\lambda + \varphi_k^{\mathsf T}P_{k-1}\varphi_k},\\
  \hat\theta_k &= \hat\theta_{k-1} + K_k\bigl(y_k - \varphi_k^{\mathsf T}\hat\theta_{k-1}\bigr),\\
  P_k &= \frac{1}{\lambda}\bigl(P_{k-1} - K_k\varphi_k^{\mathsf T}P_{k-1}\bigr),
\end{align}
其中 $\varphi_k = [\,\lvert\omega_k\rvert,\ {\Icmd}_k^{2}\,]^{\mathsf T}$，
$\lambda$ 为遗忘因子（典型取 $0.999\sim0.99999$）。
工程上有三点必须注意：
\begin{itemize}
  \item 遗忘因子越接近 $1$，参数越稳定但跟踪越慢；
  \item 必须对参数做非负钳位与变化率限幅。若参数漂到负值，
  模型会给出反物理的预测，进而导致控制发散；
  \item 只有在“真实功率反馈可信”时才启用辨识。
  裁判系统无法识别负功率，因此在回馈工况下真实功率不可信，
  此时应冻结辨识。
\end{itemize}

\paragraph{(c) 按转速误差加权分配。}
针对“走不直”，一个有效的改进是把功率预算在各轮之间按\emph{转速误差}分配，
而不是按当前消耗等比分配：
\begin{equation}\label{eq:errdist}
  P_{\max,i} = K_{\mathrm{coe}}\frac{\mathrm{err}_i}{\sum_j \mathrm{err}_j}
             + (1-K_{\mathrm{coe}})\frac{P_{\mathrm{cmd},i}}{\sum_j P_{\mathrm{cmd},j}} ,
\end{equation}
（该策略来自文献~\cite{hkustpowermodule}。）其中 $\mathrm{err}_i$ 为该轮的目标转速与实际转速之差，
$K_{\mathrm{coe}}$ 在误差总和小于下阈值时取 $0$、大于上阈值时取 $1$、中间线性插值。
直观理解是：\emph{哪个轮子“追不上”，就多给它一点功率}。
这在直线上等价于维持原状，在转弯与打滑时能显著改善姿态保持。
另一种思路~\cite{zjuwiki} 是“削转速而不是削电流”——对轮电机同比例削减\emph{目标转速}，
因为轮电机的电流是转速误差的 PID 输出，
削减目标转速不会改变运动学解算的方向性。

\paragraph{(d) 把预算分配写成小规模优化。}
最一般的形式是一个带约束的凸问题：
在 $\sum_i P_i \le \Plim$、$\lvert{\Icmd}_i\rvert\le16384$ 的约束下，
最大化某个与驾驶意图一致的效用函数（例如四轮力矩与目标力矩的加权平方误差和）。
对四变量问题可以用解析的 KKT 条件求解，计算量在 STM32F4 上是可以接受的。
这是把工程经验“升级”为可证明最优的自然方向。

% ============================================================================
\part{超级电容}

\begin{note}
本部分目前只给出总体框架与若干必要的定量关系，
用于说明超级电容在整套控制体系中处于什么位置、与 Part I 的功率控制如何分工。
详细的电路设计、控制环实现、CAN 协议与实测数据将在后续版本中补充。
\end{note}

\section{超级电容系统的总体框架}\label{sec:supercap}

\subsection{为什么需要超级电容}

底盘电机的功率需求具有强烈的\emph{脉冲性}：起步、爬坡、急停、碰撞恢复这些瞬间
需要的功率远高于匀速巡航。而规则对底盘的约束是针对\emph{平均}功率的
（通过缓冲能量来体现），电池的放电能力也受倍率限制。
这两者之间的矛盾就是超级电容存在的理由：
\textbf{用一个能量缓冲器把“瞬时峰值功率”与“平均功率”解耦}。

定量地看，一次持续 $2\ \mathrm{s}$、超出限制 $150\ \mathrm{W}$ 的爆发需要
\begin{equation}
  \Delta E = 150\ \mathrm{W}\times 2\ \mathrm{s} = 300\ \mathrm{J}
\end{equation}
的能量。若系统的能量池在 $2\ \mathrm{kJ}$ 量级，则这样的爆发可以做很多次。
这正是超电车种在比赛中能够“以低等级底盘功率打出高等级动力”的原因。

\begin{note}[与规则的关系]
除底盘功率上限之外，近年赛季还引入了整车\emph{能量池}约束
（例如对总输出能量的累计限制）。这一变化的意义在于：
功率约束只限制“瞬时能取多快”，而能量约束限制“一共能取多少”。
两者叠加之后，超电的调度策略不再只是“把峰值削平”，
还需要考虑“把有限的能量用在什么时刻”。
本部分的后续版本会结合具体规则数值展开，此处暂不深入。
\end{note}

\subsection{系统构成}

一个典型的车载超电系统由四部分组成：

\begin{enumerate}
  \item \textbf{电容组。} 由若干单体串并联构成。串联提升耐压，
  并联降低等效串联电阻（ESR）。单体之间需要均衡电路，
  否则个别单体过压会迅速失效。
  \item \textbf{双向 DC--DC 变换器。} 负责在电池母线电压与电容组电压之间双向传递能量。
  常见拓扑有串联式 Buck/Boost、并联式四开关 Buck-Boost（FSBB）与并联式两开关结构~\cite{adernal}，
  它们在效率、体积与控制难度上各有取舍。
  \item \textbf{超电控制器。} 采样母线电压电流与电容电压电流，
  运行内层控制环，并通过 CAN 与主控通信。
  \item \textbf{主控侧的接口。} 把“本周期底盘计划从电池取多少功率”这一意向告知超电控制器。
\end{enumerate}

\subsection{能量与电压的基本关系}

电容储能为
\begin{equation}\label{eq:capenergy}
  E = \tfrac{1}{2}CU^{2}.
\end{equation}
由于能量与电压的\emph{平方}成正比，可利用的能量窗口是
\begin{equation}\label{eq:usable}
  \frac{E_{\mathrm{usable}}}{E_{\max}}
  = \frac{U_{\max}^{2} - U_{\min}^{2}}{U_{\max}^{2}} .
\end{equation}
举例：某开源设计~\cite{adernal} 采用 $11$ 串 $50\ \mathrm{F}/3\ \mathrm{V}$ 单体
（单体 ESR 约 $22\ \mathrm{m}\Omega$），则
\begin{equation}
  C = \frac{50}{11} \approx 4.55\ \mathrm{F},\qquad
  U_{\max} = 33\ \mathrm{V},\qquad
  R_{\mathrm{ESR}} = 11\times 22\ \mathrm{m}\Omega \approx 242\ \mathrm{m}\Omega .
\end{equation}
充到 $30\ \mathrm{V}$ 时储能
\begin{equation}
  E = \tfrac{1}{2}\times 4.55\times 30^{2} \approx 2045\ \mathrm{J}.
\end{equation}
若允许放电到 $15\ \mathrm{V}$，则可用能量占比为
$1-(15/30)^2 = 75\%$，即约 $1534\ \mathrm{J}$——
\emph{一半的电压窗口里藏着四分之三的能量}，这是电容储能的一个基本特征。

恒功率放电时电压随时间的演化为
\begin{equation}\label{eq:capdischarge}
  U(t) = \sqrt{U_0^{2} - \frac{2Pt}{C}} ,
\end{equation}
它清楚地说明了两件事：
\begin{itemize}
  \item 电压随时间\emph{非线性}下降，越到后段下降越快；
  因此“还能撑多久”不能线性外推。
  \item 电容电压越低，同样的放电电流能提供的功率越小，
  于是可补偿的差额也越小。这正是 Part I 中
  “电容电量低于 $5\%$ 时取消额外放宽”那一条保护存在的物理原因。
\end{itemize}

\begin{note}[为什么不能把电容直接并在母线上]
若不加 DC--DC 而直接把电容组并到母线，母线电压会随电容放电大幅下降。
这会同时带来三个问题：电机在低压下性能下降；
电调的力矩电流环增益需要重新标定；
第~\ref{sec:model} 节的功率模型（其中不含母线电压）会系统性失准。
因此双向 DC--DC 不是可选项，而是必要组件——
它把“电容电压”与“母线电压”这两个量解耦开来。
\end{note}

\subsection{分层控制架构}

超电系统在控制上是\emph{两层}结构，明确这一点对调试至关重要：

\begin{table}[htbp]
  \centering
  \small
  \caption{超电系统的分层控制职责划分}
  \label{tab:layers}
  \begin{tabular}{@{}>{\raggedright\arraybackslash}p{0.14\textwidth}>{\raggedright\arraybackslash}p{0.30\textwidth}>{\raggedright\arraybackslash}p{0.26\textwidth}>{\raggedright\arraybackslash}p{0.22\textwidth}@{}}
    \toprule
    层次 & 目标 & 主要手段 & 执行位置 \\
    \midrule
    能量层（外） & 决定本周期允许底盘用多少功率，
    使缓冲能量/电容电量维持在期望区间 &
    缓冲能量闭环、电量—功率映射、充/放电模式切换 &
    超电控制器 或 主控 \\
    \addlinespace[2pt]
    力矩层（内） & 在给定功率预算下把力矩分配到各电机 &
    Part I 的功率模型与功率再分配 &
    主控 \\
    \bottomrule
  \end{tabular}
\end{table}

关于职责边界，存在两种常见做法：
\begin{itemize}
  \item \textbf{主控负责能量层。} 主控读裁判系统的缓冲能量，
  运行外环得到功率预算，再把预算通过网络告知超电控制器。
  Part I 第~\ref{sec:bufferloop} 节描述的就是这种结构。
  \item \textbf{超电控制器负责能量层。} 主控把裁判系统的数据透传给超电控制器，
  由它在本地完成能量调度与电容管理。
  这种分工的好处是能量层与执行层在同一块板子上，回路更短、延迟更低。
\end{itemize}
无论采用哪一种，有一条经验值得记住：
\begin{quote}
\textbf{当“电容还有电却仍然超功率”时，应当优先检查超电侧的功率环与硬件能力，
而不是继续调整主控的 PID。}
\end{quote}

\subsection{与功率控制的接口}

主控与超电控制器之间至少需要传递两类信息：

\begin{itemize}
  \item \textbf{下行（主控 $\to$ 超电）：本周期计划从电池取用的功率。}
  它正是 Part I 中的 $\Plim = \Pmax - u$。超电控制器据此决定
  是给电容充电（当 $\Plim$ 低于底盘的实际需求时）还是放电补足。
  \item \textbf{上行（超电 $\to$ 主控）：电容剩余电量百分比、当前母线功率、状态字。}
  电量信息用于 Part I 中的电量保护与模式切换，
  真实功率信息则可以反过来用于在线辨识电机参数（第~\ref{sec:future} 节）。
\end{itemize}

开源实现中几种简洁的接口形态见文献~\cite{hitszsupercap,pspsupercap,h7framework}，其中一种的用法是：
\begin{lstlisting}[style=cstyle, numbers=none]
input_power = max_power_limit - chassis_power_control->buffer_pid.out;
CAN_CMD_CAP(input_power);     /* 下行：告知超电本周期电池侧功率目标 */
cap_measure.cap_percent;      /* 上行：电容剩余电量百分比 */
\end{lstlisting}

典型的工作模式有两种，可以按战况切换：
\begin{description}
  \item[充电模式。] 把功率预算设为\emph{略低于}裁判限制，
  使底盘产生小幅冗余功率持续给电容充电。
  即使在交战中也能缓慢回充，是默认状态。
  \item[补电模式。] 把功率预算设为\emph{远高于}裁判限制，
  差额由电容长时间补偿。例如把 $45\ \mathrm{W}$ 的限制临时设为 $80\ \mathrm{W}$，
  则电容需持续提供 $35\ \mathrm{W}$；若电容储量约 $2000\ \mathrm{J}$，
  理论上可维持约一分钟。这一模式的代价是电容电量快速下降，
  必须在电量低于阈值时自动退回充电模式。
\end{description}

\subsection{待补充的内容}

后续版本计划在本部分补充以下内容（欢迎补充资料）：
\begin{itemize}
  \item 电容组选型与均衡电路的详细设计与实测数据；
  \item 双向 DC--DC 的拓扑对比、开关器件选型与热设计；
  \item 内层控制环（电容电流环、母线功率环）的建模与整定，
  包括“恒功率”与“恒流”两种模式的取舍；
  \item 电流采样方案（分流电阻 + 双向电流监测）的量程与精度设计，
  以及开关噪声对采样的影响；
  \item CAN 通信协议的完整报文定义与容错策略（超时、校验、状态机）；
  \item 能量调度的优化策略：在整车能量池约束下，
  如何决定“何时放电、放电多少”；
  \item 实测波形、效率曲线与故障案例。
\end{itemize}

% ============================================================================
\appendix
% 附录编号只用字母：目录的编号盒固定 2.3em 宽，"附录 A" 放进去会与标题重叠。
% 因此把编号保持为 "A"，而把"附录"二字写进标题与目录条目里。
\renewcommand{\thesection}{\Alph{section}}
\newcommand{\appsection}[1]{%
  \refstepcounter{section}%
  \phantomsection
  \section*{附录~\thesection\quad #1}%
  \addcontentsline{toc}{section}{附录~\thesection\quad #1}%
}

% ============================================================================
\appsection{矩阵求导速查与最小二乘的几何解释}\label{app:lsq}

\subsection{常用矩阵求导公式}

下表列出推导最小二乘解时用到的全部公式。
其中 $a$ 为常向量，$A$ 为常矩阵，$\theta$ 为自变量向量。

\begin{table}[htbp]
  \centering
  \small
  \caption{常用标量对向量求导公式}
  \label{tab:matrixderiv}
  \begin{tabular}{@{}lll@{}}
    \toprule
    表达式 & 导数 & 备注 \\
    \midrule
    $a^{\mathsf T}\theta$                & $a$                        & 线性型 \\
    $\theta^{\mathsf T}a$                & $a$                        & 与上式相等（标量） \\
    $\theta^{\mathsf T}\theta$           & $2\theta$                  & 取 $A=I$ \\
    $\theta^{\mathsf T}A\theta$          & $(A+A^{\mathsf T})\theta$  & 一般情形 \\
    $\theta^{\mathsf T}A\theta$          & $2A\theta$                 & $A$ 对称时 \\
    $\lVert y-X\theta\rVert^{2}$         & $-2X^{\mathsf T}y + 2X^{\mathsf T}X\theta$ & 最小二乘的核心 \\
    \bottomrule
  \end{tabular}
\end{table}

补充两条常用的标量求导法则（第~\ref{sec:firstorder} 节已用到）：
\begin{equation}
  \frac{\partial}{\partial x}\bigl[u(x)^2\bigr] = 2u\,\frac{\partial u}{\partial x},
  \qquad
  \frac{\partial}{\partial x}\bigl[u(x)v(x)\bigr]
  = \frac{\partial u}{\partial x}v + u\frac{\partial v}{\partial x}.
\end{equation}

\subsection{几何解释：最小二乘就是正交投影}

线性方程组 $X\theta = y$ 通常在 $N>3$ 时无解，因为 $y$ 一般不在 $X$ 的列空间里。
最小二乘做的事情是：在 $X$ 的列空间中找一个与 $y$ \emph{距离最近}的向量。
由第~\ref{sec:firstorder} 节得到的正交性条件 \fref{eq:orth}，
这个“最近点”正是 $y$ 在列空间上的\emph{正交投影}：
\begin{equation}
  \hat y = X\hat\theta = X(X^{\mathsf T}X)^{-1}X^{\mathsf T}y = Hy ,
\end{equation}
其中 $H = X(X^{\mathsf T}X)^{-1}X^{\mathsf T}$ 称为\emph{投影矩阵}（或称帽子矩阵），
它满足 $H^2=H$（幂等）与 $H^{\mathsf T}=H$（对称）——这正是投影算子的定义性质。
残差 $\varepsilon = y-\hat y$ 与列空间中的每一个向量都正交，
因此有著名的勾股关系
\begin{equation}
  \lVert y\rVert^{2} = \lVert \hat y\rVert^{2} + \lVert\varepsilon\rVert^{2}.
\end{equation}
“最小二乘”这个名字的来源也在这里：它最小化的正是 $y$ 到列空间的\emph{垂直距离}。

\begin{figure}[htbp]
  \centering
  \begin{tikzpicture}[scale=1.0, >=Stealth]
    % 列空间（用平行四边形示意）
    \fill[eblue!8] (0,0) -- (5.0,0.9) -- (6.6,2.9) -- (1.6,2.0) -- cycle;
    \draw[eblue!45] (0,0) -- (5.0,0.9) -- (6.6,2.9) -- (1.6,2.0) -- cycle;
    \node[eblue!70!black, font=\small] at (5.85,2.75) {$\operatorname{col}(X)$};

    \coordinate (O)  at (1.55,0.85);
    \coordinate (Yh) at (4.05,2.00);
    \coordinate (Y)  at (4.05,4.35);

    \draw[->, thick] (O) -- (Y)  node[midway, left=1pt, font=\small] {$y$};
    \draw[->, thick, eblue] (O) -- (Yh)
       node[midway, below right=0pt, font=\small, text=eblue!70!black] {$\hat y = X\hat\theta$};
    \draw[->, thick, red!70!black, dashed] (Yh) -- (Y)
       node[midway, right=2pt, font=\small, text=red!70!black] {$\varepsilon$};

    % 直角标记
    \draw[red!70!black] ($(Yh)+(0.28,0.06)$) -- ($(Yh)+(0.28,0.34)$) -- ($(Yh)+(0.0,0.28)$);
    \fill (O) circle (1.1pt);
    \node[below left, font=\small] at (O) {$0$};
  \end{tikzpicture}
  \caption{最小二乘的几何意义：$\hat y$ 是观测向量 $y$ 在回归量列空间上的正交投影，
  残差 $\varepsilon$ 与列空间正交。}
  \label{fig:projection}
\end{figure}

\subsection{从投影角度重推正规方程}

设 $\hat y = X\hat\theta$。正交条件“$\varepsilon\perp\operatorname{col}(X)$”
可以写成：对列空间中的任意向量 $Xv$（$v$ 任意），都有
\begin{equation}
  (Xv)^{\mathsf T}(y - X\hat\theta) = 0
  \;\Longrightarrow\;
  v^{\mathsf T}\bigl(X^{\mathsf T}y - X^{\mathsf T}X\hat\theta\bigr) = 0 .
\end{equation}
由于 $v$ 任意，括号内的向量必须为零向量，于是
\begin{equation}
  X^{\mathsf T}y - X^{\mathsf T}X\hat\theta = 0
  \;\Longleftrightarrow\;
  X^{\mathsf T}X\hat\theta = X^{\mathsf T}y .
\end{equation}
这与 \fref{eq:normal} 完全相同，但完全不需要求导。
两种推导各有好处：求导的版本给出了“为什么是极小值”（海森矩阵半正定），
投影的版本给出了“为什么残差正交”（几何直观）。

% ============================================================================
\appsection{单位与换算系数}

模型推导中涉及多次单位换算是出错的常见来源。下表把它们集中列出。

\begin{table}[htbp]
  \centering
  \small
  \caption{关键单位换算}
  \label{tab:units}
  \begin{tabular}{@{}lll@{}}
    \toprule
    换算 & 关系 & 数值 \\
    \midrule
    转速：rpm $\to$ rad/s      & $\Omega = \dfrac{2\pi}{60}\omega$ & $\div\,9.5493$ \\
    转速：rad/s $\to$ rpm      & $\omega = \dfrac{60}{2\pi}\Omega$ & $\times\,9.5493$ \\
    功率：扭矩 $\times$ 转速   & $P = \dfrac{\tau\omega}{9.55}$    & $\tau$ 单位 N$\cdot$m，$\omega$ 单位 rpm \\
    控制值 $\to$ 力矩电流      & $\Itau = \dfrac{20}{16384}\Icmd$  & $\times\,1.2207\times10^{-3}$ A \\
    力矩电流 $\to$ 转子力矩    & $\tau = \KT\Itau$                 & $\KT = 0.015622$ N$\cdot$m/A \\
    控制值 $\to$ 转子力矩      & $\tau = 1.9069\times10^{-5}\Icmd$ & N$\cdot$m \\
    控制值 $\to$ 机械功率系数  & $\CT = 1.99688994\times10^{-6}$   & W/rpm（每单位控制值） \\
    \bottomrule
  \end{tabular}
\end{table}

\begin{note}
在横向比较不同队伍公布的模型系数之前，必须先确认其\emph{单位约定}。
同样是写作 $k_2$ 的系数，若一方的力矩以 A 为单位、转速以 rad/s 为单位，
另一方以控制值（$\pm16384$）和 rpm 为单位，两者的数值可以相差约 $10^{6}$ 倍。
换算关系为
\begin{equation}
  k_{2}^{(\text{A},\ \mathrm{rad/s})}
  = k_{2}^{(\text{控制值},\ \mathrm{rpm})}
    \times \left(\frac{16384}{20}\right)^{2} \times 9.55 ,
\end{equation}
因此在任何对比表中都必须注明单位，否则结论没有意义。
\end{note}

% ============================================================================
\appsection{数据集与拟合脚本}

开源实现提供的\code{model\_fit} 目录包含两部分：拟合脚本与实测数据集。

\subsection{数据集}

数据集为 CSV 格式，文件名反映了激励信号的类型与参数，
包括恒定转速、斜坡、以及若干组不同频率与幅值的正弦激励：

\begin{table}[htbp]
  \centering
  \small
  \caption{数据集构成（示意）}
  \label{tab:datasets}
  \begin{tabular}{@{}lll@{}}
    \toprule
    文件 & 激励类型 & 用途 \\
    \midrule
    \code{constantspeed.csv} & 恒定转速 & 观察空载损耗随转速的变化 \\
    \code{ramp.csv}、\code{ramp2.csv} & 斜坡力矩 & 覆盖大范围力矩 \\
    \code{sin*.csv} & 正弦力矩（多种频率/幅值） & 主要拟合数据 \\
    \code{test*.csv} & 混合工况 & 交叉验证 \\
    \bottomrule
  \end{tabular}
\end{table}

如前文所述，这样的激励组合是\emph{刻意}设计的：只有同时让转速与力矩充分变化，
$k_2$ 与 $k_1$、$a$ 才可分离（第~\ref{sec:lsq} 节）。

\subsection{拟合脚本}

\begin{lstlisting}[style=matstyle]
KT = 0.3/(3591/187);
KE = (60/(2*pi*24.48))/(3591/187);

bus_voltage = 24.56;
give_toque_current = 20 * give_current / 16384;
give_toque = KT .* give_toque_current;
input_power = bus_voltage * current;

machine_power = (motor_chassis0speed_rpm .* give_toque) / 9.550;

% 损耗功率 = 实测输入功率 - 解析机械功率
Pother = input_power - machine_power;

% 线性最小二乘拟合：Pother = k1*w^2 + k2*I^2 + c
g = fittype('k1*motor_chassis0speed_rpm^2+k2*give_current.^2+c', ...
    'independent',{'motor_chassis0speed_rpm','give_current'}, ...
    'dependent','Pother','coefficients',{'k1','k2','c'});
myfit = fit([motor_chassis0speed_rpm, give_current], Pother, g);

% 拟合结果示例
Pre_Pother = 2.11e-07*motor_chassis0speed_rpm.^2 ...
           + 9.805e-08*give_current.^2 + 2.138;
predicte_power = machine_power + Pre_Pother;

figure; plot(input_power); hold on; plot(predicte_power);
legend('real power','predicte power');
\end{lstlisting}

\code{fittype} 中的表达式与本文 \fref{eq:regmodel} 完全一致，
只是变量名不同：\code{motor\_chassis0speed\_rpm} 对应 $\omega$，
\code{give\_current} 对应 $\Icmd$，\code{Pother} 对应 $y$。
\code{fit} 内部使用的是正交分解而非显式求逆，
因此第~\ref{sec:lsq} 节讨论的数值条件问题在 MATLAB 中通常不会暴露——
但这也意味着\textbf{直接照抄公式到单片机上可能得到不准确的结果}，
需要按第~\ref{sec:lsq} 节的建议先做无量纲化。

\subsection{自己实现正规方程}

若需要在 MATLAB 之外（例如 Python 或嵌入式离线工具）复现拟合，
最小实现只需要构造 $X$ 与 $y$，然后求解 \fref{eq:normal}：

\begin{lstlisting}[style=cstyle, numbers=none]
# X 的第 i 行为 [omega_i^2, Icmd_i^2, 1]，y_i 为损耗功率
import numpy as np
X = np.column_stack([omega**2, Icmd**2, np.ones_like(omega)])
theta, residuals, rank, sv = np.linalg.lstsq(X, y, rcond=None)
# theta = [k2, k1, a]
print("cond(X) =", np.linalg.cond(X))
\end{lstlisting}

顺手打印条件数是一个好习惯：若 $\kappa(X)$ 大于 $10^{6}$，
说明数据存在严重共线性，此时参数虽然可以算出来，但不宜相信。

% ============================================================================
\appsection{中英文术语对照}

\begin{table}[htbp]
  \centering
  \small
  \caption{术语对照表}
  \label{tab:glossary}
  \begin{tabular}{@{}lll@{}}
    \toprule
    中文 & 英文 / 代码标识 & 说明 \\
    \midrule
    缓冲能量           & buffer energy / \code{chassis\_power\_buffer} & 裁判系统给出的功率余量 \\
    缓冲能量闭环       & buffer power loop            & 外环，PI 控制 \\
    功率上限           & power limit / \code{max\_power\_limit} & 裁判系统给定 \\
    功率预算           & power budget / \code{chassis\_max\_power} & 本周期允许的取用功率 \\
    力矩电流控制值     & give current / \code{give\_current} & 发送给电调的整型控制量 \\
    力矩电流反馈值     & given current / \code{given\_current} & 电调回传的控制量 \\
    力矩电流           & torque current / $I_\tau$    & 与转矩成正比的电流分量 \\
    力矩电流系数       & \code{toque\_coefficient} / $\CT$ & 发送值 $\times$ 转速 $\to$ 功率 \\
    空载损耗           & no-load loss / $P_0(\omega)$ & 零力矩时的输入功率 \\
    功率再分配         & power redistribution         & 本文采用的分配策略 \\
    功率缩放系数       & power scale / \code{power\_scale} & $k=\Plim/\sum P_i$ \\
    恒功率包络         & constant-power envelope      & 允许控制量关于转速的曲线 \\
    回馈制动           & regenerative braking         & 力矩与转速反向，功率为负 \\
    反电动势           & back electromotive force     & 减速时产生回馈功率的物理机制 \\
    磁场定向控制       & field-oriented control (FOC) & 将电流分解为 $d$、$q$ 轴 \\
    递推最小二乘       & recursive least squares (RLS) & 在线参数辨识 \\
    遗忘因子           & forgetting factor $\lambda$  & RLS 对新数据的权重 \\
    等效串联电阻       & equivalent series resistance (ESR) & 电容组的关键参数 \\
    填谷削峰           & peak shaving / valley filling & 超电的基本工作方式 \\
    电池管理系统       & battery management system (BMS) & — \\
    \bottomrule
  \end{tabular}
\end{table}

% ============================================================================
\printbibliography[heading=bibintoc, title=\ebibname]

\end{document}
